\section{Scope and proof contract}

The goal is kernel-checked rigidity at perimetral densities $1/3$, $1/4$,
and $1/5$, supported by reusable lemmas about curves and moving polygons.
This is an ongoing blueprint, not a completed formalization. Geometric
rigidity at $1/3$ and $1/4$ is proved in Lean; the $1/5$ target remains
unproved. The scope of the checked theorems is
compact strictly convex planar bodies with a specified simple, $C^1$,
unit-speed boundary parametrization. No central symmetry is assumed.

Perimetral density $\sigma$ is a boundary-length fraction, distinct from
submerged area fraction $\rho$. We first target the geometric constant-chord
formulation. Arc-length normalization of regular boundary data is proved.
The connection to physical equilibrium remains a separate proof obligation.

Only kernel-checked declarations may carry \texttt{leanok}. Defining a
target proposition is not proving it. Each statement-only link ends in
\texttt{Statement} and names a definition of type \texttt{Prop}, not a proof
or an axiom. Such nodes remain unmarked until a separately named theorem
proves the specified proposition. Mathlib quality is a design goal, not a claim
of upstream review or acceptance.

\section{General definitions and first lemmas}

\begin{definition}[Constant chord]
\label{def:chord}
\lean{UlamFloating.HasConstantChord}
\leanok
For $\gamma:\R\to E$ into a pseudometric space, constant chords of increment
$h$ and length $\ell$ mean $d(\gamma(t+h),\gamma(t))=\ell$ for every $t$.
\end{definition}

\begin{lemma}[Basic chord API]
\label{lem:chord-api}
\lean{UlamFloating.HasConstantChord.nonneg,UlamFloating.HasConstantChord.neg,UlamFloating.HasConstantChord.shift,UlamFloating.HasConstantChord.complement}
\leanok
\uses{def:chord}
The length is nonnegative. Negating the increment or translating the
parameter preserves the chord law. If $\gamma$ is $P$-periodic, replacing
$h$ by $P-h$ also preserves it.
\end{lemma}
\begin{proof}
\leanok
Use distance symmetry and nonnegativity. Evaluate the law at $t-h$ for
reversal, then use $\gamma(t+P-h)=\gamma(t-h)$ for the complement.
\end{proof}

\begin{definition}[Boundary data]
\label{def:boundary}
\lean{UlamFloating.Plane,UlamFloating.IsUnitSpeedBoundary}
\leanok
Use Euclidean $\R^2$, with its Euclidean norm. Boundary data for $K$ consist
of $P>0$ and $\gamma:\R\to\R^2$, where $K$ is compact, strictly convex,
and has nonempty interior; $\gamma$ is $C^1$ and $P$-periodic; and
\[
 \norm{\gamma'(t)}=1,\quad \gamma(\R)=\partial K,\quad
 \gamma(s)=\gamma(t)\ \Longleftrightarrow\
 \exists k\in\Z,\ s=t+kP.
\]
The last condition excludes multiple traversals.
\end{definition}

\begin{definition}[Perimetral property and rigidity target]
\label{def:perimetral}
\lean{UlamFloating.HasPerimetralChordProperty,UlamFloating.IsDisk,UlamFloating.PerimetralRigidity}
\leanok
\uses{def:boundary,def:chord}
Let $\mathcal P_\sigma(K)$ mean $0<\sigma<1$ and there exist boundary data
and $\ell>0$ with constant chords of increment $\sigma P$. A disk is a
closed Euclidean ball of positive radius. Define the proposition
\[
 \mathcal R(\sigma):\quad
 \forall K,\ \mathcal P_\sigma(K)\Longrightarrow
 \exists c\in\R^2\ \exists r>0,\ K=\overline B(c,r).
\]
\end{definition}

\begin{lemma}[Complementary densities]
\label{lem:complement}
\lean{UlamFloating.HasPerimetralChordProperty.complement,UlamFloating.PerimetralRigidity.complement}
\leanok
\uses{def:perimetral,lem:chord-api}
$\mathcal P_\sigma(K)$ implies $\mathcal P_{1-\sigma}(K)$, and
$\mathcal R(\sigma)$ implies $\mathcal R(1-\sigma)$.
\end{lemma}
\begin{proof}
\leanok
Apply the chord complement with $(1-\sigma)P=P-\sigma P$.
Apply it twice for the rigidity implication.
\end{proof}

\begin{definition}[Regular boundary and traversal length]
\label{def:regular-boundary}
\lean{UlamFloating.IsRegularBoundary,UlamFloating.traversalLength}
\leanok
\uses{def:boundary}
A regular boundary has the compactness, convexity, simplicity, periodicity,
and range conditions of boundary data, with nonzero derivative replacing
unit speed. For a period $Q$, its traversal length is
$\int_0^Q\norm{\varphi'(t)}\,dt$.
\end{definition}

\begin{lemma}[Variation bounds for one-Lipschitz curves]
\label{lem:lipschitz-variation}
\lean{UlamFloating.eVariationOn_le_of_one_lipschitz,UlamFloating.variation_bounds_of_one_lipschitz}
\leanok
A one-Lipschitz curve into any pseudometric space has variation on
$[a,b]$ at most $b-a$ for $a\leq b$. Its endpoint distance is a lower
bound for that variation.
\end{lemma}
\begin{proof}
\leanok
Apply the Lipschitz composition bound for variation to the identity
function. Its variation is bounded by the difference of the endpoints.
The two-point partition supplies the lower bound.
\end{proof}

\begin{lemma}[Derivative of cumulative variation]
\label{lem:variation-derivative}
\lean{UlamFloating.slope_variation_bounds,UlamFloating.hasDerivAt_variation_of_unit_speed}
\leanok
\uses{lem:lipschitz-variation}
For a one-Lipschitz curve, difference quotients of signed cumulative
variation lie between the norm of the curve's difference quotients and
one. For a differentiable unit-speed curve they therefore converge to
one. Continuity of the derivative is unnecessary.
\end{lemma}
\begin{proof}
\leanok
Additivity of signed variation identifies its increments with interval
variation. Divide the preceding bounds by the parameter increment, treating
both orders of endpoints. Use the derivative of the curve and the squeeze
theorem to obtain derivative one for cumulative variation.
\end{proof}

\begin{proposition}[Metric length of a unit-speed curve]
\label{prop:unit-speed-length}
\lean{UlamFloating.variation_eq_sub_of_unit_speed,UlamFloating.eVariationOn_eq_of_unit_speed}
\leanok
\uses{lem:variation-derivative}
For every differentiable unit-speed curve in a real normed space, signed
variation from $a$ to $b$ equals $b-a$. When $a\leq b$, its metric total
variation on $[a,b]$ is exactly $b-a$.
\end{proposition}
\begin{proof}
\leanok
Subtract the parameter from cumulative variation. Its derivative vanishes,
so the mean value theorem makes the difference constant. Variation at the
base point is zero. The one-Lipschitz bound guarantees finite variation,
allowing conversion back from real-valued signed variation.
\end{proof}

\begin{proposition}[Length of a unit-speed boundary traversal]
\label{prop:boundary-length}
\lean{UlamFloating.IsUnitSpeedBoundary.traversalLength_eq,UlamFloating.unit_speed_boundary_length}
\leanok
\uses{def:boundary,def:regular-boundary,prop:unit-speed-length}
For unit-speed boundary data, both the integral of speed and metric total
variation over $[0,P]$ equal $P$.
\end{proposition}
\begin{proof}
\leanok
Integrate the constant speed one, and apply the unit-speed metric length
theorem to the positive traversal interval.
\end{proof}

\begin{lemma}[Integral shifts by integer periods]
\label{lem:integral-period-shift}
\lean{UlamFloating.integral_add_int_mul_period}
\leanok
For a continuous $Q$-periodic real function $v$ and every integer $k$,
\[
 \int_0^{t+kQ}v=\int_0^t v+k\int_0^Qv.
\]
\end{lemma}
\begin{proof}
\leanok
Split the integral after shifting the upper endpoint by the period $kQ$,
and use the integer-multiple formula for integrals of periodic functions.
\end{proof}

\begin{lemma}[Integral coordinates for positive periodic speed]
\label{lem:integral-coordinates}
\lean{UlamFloating.strictMono_integral_of_pos,UlamFloating.surjective_integral_of_pos_periodic,UlamFloating.exists_integral_orderIso}
\leanok
\uses{lem:integral-period-shift}
For a positive continuous $Q$-periodic function $v$, with $Q>0$,
$S(t)=\int_0^t v$ is an increasing $C^1$ order isomorphism of $\R$.
Its inverse is $C^1$, with derivative $1/v(S^{-1}(t))$.
\end{lemma}
\begin{proof}
\leanok
The fundamental theorem gives $S'=v>0$, hence strict monotonicity.
If $P=\int_0^Qv>0$, then $S(kQ)=kP$. Integer multiples bracket every
real value, so continuity gives surjectivity. Bundle this map as an
order isomorphism and apply the inverse derivative and regularity theorems.
\end{proof}

\begin{proposition}[Unit-speed reparametrization in a real normed space]
\label{prop:regular-curve-reparam}
\lean{UlamFloating.exists_unitSpeed_reparam}
\leanok
\uses{lem:integral-coordinates,lem:integral-period-shift,prop:unit-speed-length}
Every regular $C^1$ periodic curve in a real normed space admits an
increasing $C^1$ change of parameter $e$, with $C^1$ inverse and $e(0)=0$,
such that $\varphi\circ e$ has unit speed. For the new positive period
$P=\int_0^Q\norm{\varphi'}$, one has
$e(t+kP)=e(t)+kQ$ for every integer $k$. The original curve's metric
total variation on $[0,Q]$ equals $P$.
\end{proposition}
\begin{proof}
\leanok
Apply the integral-coordinate theorem to $v=\norm{\varphi'}$ and put
$e=S^{-1}$. Differentiate the composition to obtain unit speed.
The integer shift law for $S$ gives the corresponding law for its inverse.
Increasing reparametrization preserves metric variation; $S$ maps
$[0,Q]$ onto $[0,P]$, where the unit-speed length theorem applies.
\end{proof}

\begin{proposition}[Arc-length normalization of boundary data]
\label{prop:arclength}
\lean{UlamFloating.arc_length_normalization}
\leanok
\uses{def:boundary,def:regular-boundary,prop:regular-curve-reparam}
Given a regular boundary $(K,Q,\varphi)$, there exist $P$, a unit-speed
boundary $\gamma$, and an increasing $C^1$ order isomorphism $e:\R\to\R$
with $C^1$ inverse, such that $e(0)=0$, $e(t+P)=e(t)+Q$, and
$\gamma(t)=\varphi(e(t))$. The number $P$ equals both the integral of speed
and the metric total variation of $\varphi$ on $[0,Q]$.
\end{proposition}
\begin{proof}
\leanok
Apply the normed-space reparametrization theorem. The integer shift law
preserves once-traversed simplicity in both directions, while the order
isomorphism's surjectivity preserves the range. The underlying body is
unchanged, so its compactness and strict convexity also hold for the
normalized boundary data. Metric length is expressed by mathlib's
\texttt{eVariationOn}. The input is a regular embedded periodic curve;
a separate manifold-only boundary interface is not assumed here.
\end{proof}

\section{Rational sampling and local kinematics}

\begin{definition}[Sampled vertices]
\label{def:sampling}
\lean{UlamFloating.sampledVertex}
\leanok
Set $v_i(t)=\gamma(t+ih)$ for natural-number indices $i$.
\end{definition}

\begin{lemma}[Closure and equal edges]
\label{lem:sampling}
\lean{UlamFloating.sampledVertex_zero,UlamFloating.sampledVertex_succ,UlamFloating.sampledVertex_add_period,UlamFloating.HasConstantChord.sampled_edge}
\leanok
\uses{def:sampling,def:chord}
For $n\neq0$ and $h=P/n$, a $P$-periodic curve satisfies
$v_{i+n}(t)=v_i(t)$. The chord law gives
$d(v_{i+1}(t),v_i(t))=\ell$. Also $v_0(t)=\gamma(t)$ and
$v_{i+1}(t)=v_i(t+h)$.
\end{lemma}
\begin{proof}
\leanok
Expand the parameter, use $n(P/n)=P$, and apply periodicity.
Evaluate the chord law at $t+ih$ for the edge equality.
\end{proof}

\begin{definition}[Carousel]
\label{def:carousel}
\lean{UlamFloating.IsCarousel}
\leanok
For a permutation $s$ of an index type, a carousel consists of regular
$C^1$ curves $p_i$ in a real inner product space, a constant $\ell>0$, and
the laws
\[
 \norm{p_{s(i)}-p_i}=\ell,\qquad
 p_i'+p_{s(i)}'=a_i(t)(p_{s(i)}-p_i)
\]
for some real scalar $a_i(t)$ at every index and time. The velocity sum is
twice the midpoint velocity. One cyclic orbit, convexity, and the fact that
all vertices trace the same boundary are additional application hypotheses.
\end{definition}

\begin{lemma}[Rigid-edge derivative]
\label{lem:orthogonal}
\lean{UlamFloating.inner_edge_velocity_eq_zero}
\leanok
If $e=q-p$ has constant norm, and $u=p'$, $v=q'$, then $\ip e{v-u}=0$.
\end{lemma}
\begin{proof}
\leanok
Differentiate $\ip ee=\ell^2$ and use symmetry of the real inner product.
\end{proof}

\begin{lemma}[Equal speeds and reflection]
\label{lem:reflection}
\lean{UlamFloating.norm_eq_of_carousel_laws,UlamFloating.velocity_eq_reflection}
\leanok
If $\ip e{v-u}=0$ and $u+v=ae$, then $\norm u=\norm v$.
For $e\neq0$, $v=(2\ip eu/\norm e^2)e-u$.
\end{lemma}
\begin{proof}
\leanok
Pair the velocity-sum equation with $v-u$ to compare squared norms.
Pair it with $e$ to solve for $a$. No angle lifts are required.
\end{proof}

\begin{definition}[Finite polygon geometry]
\label{def:polygon-geometry}
\lean{UlamFloating.planeAreaForm,UlamFloating.boundarySamples,UlamFloating.IsStrictlyConvexPolygon,UlamFloating.IsRegularPolygon}
\leanok
\uses{def:sampling}
Finite samples use indices in $\operatorname{Fin}(n)$ and successor
\texttt{finRotate}. A strictly convex polygon has distinct vertices, $n\geq3$,
and an orientation for which every nonincident vertex lies strictly to the
left of every oriented edge. A regular polygon has a center $c$, radius
$r>0$, and the full distance relations
\[
 d(p_i,c)=r,\qquad d(p_i,p_j)^2=
 2r^2\bigl(1-\cos(2\pi(i-j)/n)\bigr).
\]
The full pairwise formula specifies cyclic labeling and excludes star order.
The area form is mathlib's oriented Euclidean area form, specialized to the plane.
\end{definition}

\begin{lemma}[Finite cyclic sampling and differentiation]
\label{lem:finite-cycle}
\lean{UlamFloating.boundarySamples_rotate,UlamFloating.HasConstantChord.boundarySamples_edge,UlamFloating.hasDerivAt_boundarySamples,UlamFloating.deriv_boundarySamples,UlamFloating.deriv_boundarySamples_rotate}
\leanok
\uses{lem:sampling,def:polygon-geometry}
Finite samples satisfy $p_{i+1}(t)=p_i(t+P/n)$, including the last-to-first
wrap. The constant-chord law gives every cyclic edge length.
Differentiating a sample translates the original derivative by its
parameter offset; differentiating the cyclic shift gives the next tangent.
\end{lemma}
\begin{proof}
\leanok
Split into the last index and all other indices. Use periodicity for the
wrap and the natural-number sampling identity elsewhere. Apply the chain
rule to the parameter translations.
\end{proof}

\begin{lemma}[Boundary traversal and distinct samples]
\label{lem:boundary-embedding}
\lean{UlamFloating.IsUnitSpeedBoundary.deriv_ne_zero,UlamFloating.IsUnitSpeedBoundary.mem_frontier,UlamFloating.IsUnitSpeedBoundary.mem,UlamFloating.IsUnitSpeedBoundary.injOn_Ico,UlamFloating.IsUnitSpeedBoundary.boundarySamples_injective}
\leanok
\uses{def:boundary,def:polygon-geometry}
Every boundary value lies in the body and its frontier, and every tangent
is nonzero. The boundary is injective on each half-open interval $[a,a+P)$.
Consequently the finite samples at distinct indices are distinct points.
\end{lemma}
\begin{proof}
\leanok
Equal boundary values have parameters differing by an integer multiple
of $P$. Two parameters in the stated interval force that integer strictly
between $-1$ and $1$, hence zero. Sample offsets lie in $[0,P)$.
\end{proof}

\begin{lemma}[Strict supporting functionals and tangents]
\label{lem:support}
\lean{UlamFloating.exists_strict_support,UlamFloating.exists_strict_support_tangent}
\leanok
\uses{def:boundary}
For a strictly convex subset of a real normed space and a noninterior
point $p$ in it, there is a continuous linear functional $f$ with
$f(x)<f(p)$ for every other point $x$ of the set. If a differentiable
curve stays in the set and passes through $p$, its tangent $u$ satisfies
$f(u)=0$.
\end{lemma}
\begin{proof}
\leanok
Separate $p$ from the open convex interior by Hahn--Banach. The midpoint
of $p$ and any other point lies in the interior, yielding strict separation
on the whole set. Along the curve, the supporting functional has a maximum
at $p$; Fermat's theorem gives zero derivative.
\end{proof}

\begin{lemma}[Planar decomposition and tangent sides]
\label{lem:planar-decomposition}
\lean{UlamFloating.norm_sq_smul_eq_inner_add_area,UlamFloating.areaForm_mul_pos_of_support}
\leanok
\uses{def:polygon-geometry}
For the oriented area form $\omega$ and right-angle rotation $J$,
\[
 \norm u^2 x=\ip ux\,u+\omega(u,x)Ju.
\]
If $\norm u=1$, $f(u)=0$, and $f(x),f(y)<0$, then
$\omega(u,x)\omega(u,y)>0$.
\end{lemma}
\begin{proof}
\leanok
Pair with an arbitrary vector and use the two-dimensional Gram identity.
Applying $f$ to the decomposition gives $f(x)=\omega(u,x)f(Ju)$;
multiply the corresponding equalities for $x$ and $y$.
\end{proof}

\begin{lemma}[Consistent directed tangent orientation]
\label{lem:tangent-orientation}
\lean{UlamFloating.mul_pos_of_continuous_ne_zero,UlamFloating.IsUnitSpeedBoundary.areaForm_mul_pos,UlamFloating.IsUnitSpeedBoundary.tangent_area_ne_zero,UlamFloating.IsUnitSpeedBoundary.tangent_area_mul_pos}
\leanok
\uses{lem:support,lem:planar-decomposition,lem:boundary-embedding}
Fix an interior point $c$ and an orientation. The continuous function
$D(t)=\omega(\gamma'(t),c-\gamma(t))$ never vanishes and satisfies
$D(s)D(t)>0$ for all $s,t$. For any other point $x$ of the body,
\[
 \omega(\gamma'(t),x-\gamma(t))D(t)>0.
\]
The constant-sign argument is proved for any continuous nonvanishing
real function on a preconnected space.
\end{lemma}
\begin{proof}
\leanok
Use a strict supporting functional and planar decomposition for the
pointwise inequalities. Continuity follows from the $C^1$ hypothesis.
A change of sign would give a zero by the intermediate value theorem.
\end{proof}

\begin{lemma}[Distinct directed tangents]
\label{lem:distinct-tangents}
\lean{UlamFloating.IsUnitSpeedBoundary.deriv_ne_of_ne}
\leanok
\uses{lem:tangent-orientation}
If $\gamma(s)\neq\gamma(t)$, then $\gamma'(s)\neq\gamma'(t)$.
\end{lemma}
\begin{proof}
\leanok
If both tangents equal $u$, put $A=\omega(u,\gamma(t)-\gamma(s))$.
The supporting-side inequalities give $A D(s)>0$ and $-A D(t)>0$.
Multiplying contradicts $D(s)D(t)>0$.
\end{proof}

\begin{lemma}[Midpoint law from convex boundary geometry]
\label{lem:midpoint-extraction}
\lean{UlamFloating.exists_smul_of_inner_eq_zero,UlamFloating.midpoint_parallel_of_equal_norms,UlamFloating.IsUnitSpeedBoundary.chord_midpoint_parallel}
\leanok
\uses{lem:orthogonal,lem:planar-decomposition,lem:distinct-tangents}
In the plane, if $e\neq0$, $u\neq v$, $\norm u=\norm v$, and
$\ip e{v-u}=0$, then $u+v$ is parallel to $e$.
Thus every constant positive chord of the strictly convex unit-speed
boundary has edge-parallel midpoint velocity.
\end{lemma}
\begin{proof}
\leanok
Both $e$ and $u+v$ are perpendicular to the nonzero vector $v-u$.
The planar Gram identity makes them parallel. Positive chord length and
the distinct-tangent theorem supply the two nonzero hypotheses.
\end{proof}

\begin{proposition}[Carousel kinematic extraction]
\label{prop:kinematic-extraction}
\lean{UlamFloating.IsUnitSpeedBoundary.boundarySamples_isCarousel}
\leanok
\uses{lem:finite-cycle,lem:boundary-embedding,lem:midpoint-extraction,def:carousel}
Every rational sampling of a strictly convex unit-speed boundary with
constant positive chords is a carousel: it has regular $C^1$ vertices,
equal constant edge lengths, and edge-parallel midpoint velocities.
\end{proposition}
\begin{proof}
\leanok
Combine the finite-sampling derivative formulas, unit speed, cyclic
edge equalities, and the midpoint law at each translated parameter.
\end{proof}

\begin{lemma}[Distinct frontier points are noncollinear]
\label{lem:frontier-triangle}
\lean{UlamFloating.areaForm_ne_zero_of_mem_frontier}
\leanok
\uses{lem:support,lem:planar-decomposition}
Three distinct frontier points of a closed strictly convex planar set
have nonzero oriented area.
\end{lemma}
\begin{proof}
\leanok
If $c-a=k(b-a)$, strict support at $a$ gives $k>0$ and strict support at
$b$ gives $k<1$. Thus $c=(1-k)a+kb$ lies in the interior by strict
convexity, contradicting its frontier membership.
\end{proof}

\begin{lemma}[Positive sign from a boundary derivative]
\label{lem:interval-sign}
\lean{UlamFloating.pos_on_Ioo_of_hasDerivAt_pos}
\leanok
Let $f$ be continuous and nonzero on $(a,b)$, with $f(a)=0$ and a
positive derivative at $a$. Then $f(x)>0$ for every $x\in(a,b)$.
\end{lemma}
\begin{proof}
\leanok
The right difference quotients converge to a positive number, so $f$
is positive at some point in $(a,b)$. A negative value elsewhere would
force a zero by the intermediate value theorem.
\end{proof}

\begin{lemma}[Positive boundary orientation]
\label{lem:positive-orientation}
\lean{UlamFloating.IsUnitSpeedBoundary.exists_positive_orientation}
\leanok
\uses{lem:tangent-orientation}
There is an orientation $\omega$ such that
$\omega(\gamma'(t),x-\gamma(t))>0$ for every $t$ and every
$x\in K\setminus\{\gamma(t)\}$.
\end{lemma}
\begin{proof}
\leanok
Choose the sign of the orientation using one tangent and an interior
point. The previously proved sign products propagate positivity to
every tangent and every other point of the body.
\end{proof}

\begin{lemma}[Oriented area respects traversal order]
\label{lem:ordered-triangle}
\lean{UlamFloating.IsUnitSpeedBoundary.areaForm_pos_of_ordered,UlamFloating.areaForm_sub_cyclic}
\leanok
\uses{lem:frontier-triangle,lem:interval-sign,lem:positive-orientation,lem:boundary-embedding}
In the positive boundary orientation, $a<b<c<a+P$ implies
$\omega(\gamma(b)-\gamma(a),\gamma(c)-\gamma(a))>0$.
Cyclically rotating the three vertices preserves this area.
\end{lemma}
\begin{proof}
\leanok
Fix $a,c$ and vary $b$. The area is continuous and nonzero on $(a,c)$
by injectivity and noncollinearity. It vanishes at $a$ and its derivative
there is strictly positive by the tangent-side inequality.
Cyclic invariance follows by bilinearity and skew symmetry.
\end{proof}

\begin{proposition}[Strict convexity of rational boundary samples]
\label{prop:sample-convexity}
\lean{UlamFloating.IsUnitSpeedBoundary.boundarySamples_strictlyConvex}
\leanok
\uses{lem:ordered-triangle,lem:boundary-embedding,def:polygon-geometry}
For every $n\geq3$ and every $t$, the rational boundary samples form
a strictly convex polygon in their specified cyclic order. No constant
chord-length assumption is required for this conclusion.
\end{proposition}
\begin{proof}
\leanok
For each directed edge, order its endpoints and any nonincident vertex
within one traversal. Apply the ordered-triangle inequality and cyclic
area invariance, including the edge from the last sample to the first.
\end{proof}

\begin{proposition}[Full convex carousel extraction]
\label{prop:extraction}
\lean{UlamFloating.carousel_extraction}
\leanok
\uses{prop:kinematic-extraction,prop:sample-convexity,lem:finite-cycle}
For $n\geq3$, the samples from $\mathcal P_{1/n}(K)$ descend to
$\Z/n\Z$ and form a strictly convex equilateral carousel.
\end{proposition}
\begin{proof}
\leanok
Combine kinematic extraction, strict convexity of the sampled polygon,
and the finite cyclic time-shift identity.
\end{proof}

\begin{lemma}[Intrinsic second carousel law]
\label{lem:second-carousel}
\lean{UlamFloating.velocity_sub_eq_rotation,UlamFloating.abs_angular_speed_le,UlamFloating.IsCarousel.edge_rotation}
\leanok
\uses{def:carousel,lem:reflection}
For an oriented planar carousel edge $e=q-p$ of length $\ell$ and endpoint
velocities $u,v$,
\[
 v-u=\frac{2\omega(u,e)}{\ell^2}Je.
\]
When $\norm u=1$, the absolute value of this angular coefficient is at
most $2/\ell$.
\end{lemma}
\begin{proof}
\leanok
The midpoint law gives $\omega(e,u+v)=0$. Orthogonality of edge and
relative velocity, together with the oriented orthogonal decomposition,
gives the displayed formula. The area-form bound
$|\omega(u,e)|\leq\norm u\norm e$ gives the estimate.
\end{proof}

\begin{proposition}[Positive rotation of constant boundary chords]
\label{prop:boundary-chord-rotation}
\lean{UlamFloating.IsUnitSpeedBoundary.chord_rotation}
\leanok
\uses{lem:second-carousel,lem:positive-orientation,prop:kinematic-extraction}
For a strictly convex unit-speed boundary with constant positive chords
of length $\ell$, there is one orientation in which the chord vector
$e(t)=\gamma(t+h)-\gamma(t)$ satisfies $e'=aJe$, where
$0<a(t)\leq2/\ell$ for every $t$.
\end{proposition}
\begin{proof}
\leanok
Choose the positive tangent orientation. The other chord endpoint lies
strictly to the left of the tangent, making $\omega(\gamma',e)>0$.
Differentiate the chord vector and use the intrinsic second carousel law.
This establishes rotation without choosing a global angle lift.
\end{proof}

\begin{lemma}[Reconstruction from angular speed]
\label{lem:rotation-reconstruction}
\lean{UlamFloating.rotation_solution_eq,UlamFloating.rotation_solution_eq_integral}
\leanok
In an oriented Euclidean plane, suppose $e'=aJe$ and $\theta'=a$ with
$\theta(0)=0$. Then
\[
 e(t)=\cos(\theta(t))e(0)+\sin(\theta(t))Je(0).
\]
For continuous $a$, one may take $\theta(t)=\int_0^t a$.
\end{lemma}
\begin{proof}
\leanok
Differentiate $\cos(\theta)e-\sin(\theta)Je$. The derivative vanishes
because $J^2=-1$, so the expression equals $e(0)$. Rotate back using
$\cos^2\theta+\sin^2\theta=1$. The integral version follows by the
fundamental theorem of calculus. No argument-function branch is used.
\end{proof}

\begin{proposition}[Global real coordinates for the chord angle]
\label{prop:global-chord-angle}
\lean{UlamFloating.IsUnitSpeedBoundary.chord_angle_coordinates}
\leanok
\uses{prop:boundary-chord-rotation,lem:rotation-reconstruction,lem:integral-coordinates,lem:integral-period-shift}
Constant positive boundary chords admit an orientation and a strictly
increasing $C^1$ order isomorphism $\theta:\R\to\R$, with $C^1$ inverse,
$\theta(0)=0$, and $0<\theta'\leq2/\ell$. The chord vector is reconstructed
by the preceding formula, and
\[
 \theta(t+kP)=\theta(t)+k\theta(P)\qquad(k\in\Z).
\]
\end{proposition}
\begin{proof}
\leanok
The angular speed from the chord-rotation law is continuous, positive,
and $P$-periodic. Apply integral coordinates, the reconstruction formula,
and the integer-period integral identity.
\end{proof}

\begin{lemma}[Nonparallel chords from cyclic convex order]
\label{lem:nonparallel-chords}
\lean{UlamFloating.IsUnitSpeedBoundary.chord_not_parallel_of_interlaced,UlamFloating.IsUnitSpeedBoundary.chord_not_parallel_of_adjacent}
\leanok
\uses{lem:ordered-triangle}
For $a<b<c<d<a+P$, the vectors $\gamma(d)-\gamma(b)$ and
$\gamma(c)-\gamma(a)$ are not parallel. For $a<b<c<a+P$, the consecutive
edge vectors are not parallel either.
\end{lemma}
\begin{proof}
\leanok
In the positive orientation, the two endpoints of an interlaced diagonal
lie on opposite sides of the other diagonal. A parallel displacement
would leave their signed side coordinates equal. For consecutive edges,
parallelism would make the ordered triangle have zero area.
\end{proof}

\begin{lemma}[No intermediate return of a directed chord vector]
\label{lem:chord-nonreturn}
\lean{UlamFloating.IsUnitSpeedBoundary.chord_ne_of_interlaced,UlamFloating.IsUnitSpeedBoundary.chord_ne_of_disjoint,UlamFloating.IsUnitSpeedBoundary.chord_ne_of_adjacent,UlamFloating.IsUnitSpeedBoundary.chord_vector_ne_initial}
\leanok
\uses{lem:nonparallel-chords,lem:ordered-triangle}
If $0<h$ and $2h<P$, then for $0<t<P$,
\[
 \gamma(t+h)-\gamma(t)\ne\gamma(h)-\gamma(0).
\]
Constant chord length is not required.
\end{lemma}
\begin{proof}
\leanok
Split according to whether the two chord intervals overlap, meet at an
endpoint, are disjoint, or cross the period endpoint. Interlaced and
adjacent cases use nonparallelism. In the disjoint case, equal directed
vectors contradict the positive side inequalities for the cyclic edges.
Periodicity reduces the wraparound cases to the same arguments.
\end{proof}

\begin{lemma}[Integral number of turns]
\label{lem:rotation-integer-turns}
\lean{UlamFloating.rotation_angle_eq_int_mul}
\leanok
If rotation through a real angle $\theta$ fixes a nonzero vector in an
oriented Euclidean plane, then $\theta=2k\pi$ for some integer $k$.
\end{lemma}
\begin{proof}
\leanok
Take the inner product with the fixed vector. Orthogonality to its
right-angle rotation gives $\cos\theta=1$. Use the exact characterization
of the real angles with cosine one.
\end{proof}

\begin{lemma}[One turn from first return]
\label{lem:angle-first-return}
\lean{UlamFloating.angle_eq_two_pi_of_no_return}
\leanok
\uses{lem:rotation-integer-turns}
Suppose a nonzero rotating vector has increasing real angular coordinates
given by an order isomorphism, starts at angle zero, and first returns at
$P>0$. Its angular increment at $P$ is exactly $2\pi$.
\end{lemma}
\begin{proof}
\leanok
The endpoint increment is a positive integer multiple of $2\pi$.
If it exceeded $2\pi$, the inverse coordinate at $2\pi$ would give an
intermediate return. Thus the integer is one.
\end{proof}

\begin{proposition}[Exactly one turn of constant boundary chords]
\label{prop:chord-one-turn}
\lean{UlamFloating.IsUnitSpeedBoundary.chord_angle_coordinates_one_turn}
\leanok
\uses{prop:global-chord-angle,lem:chord-nonreturn,lem:angle-first-return}
When $0<h$ and $2h<P$, the constructed chord angle satisfies
$\theta(P)=2\pi$ and $\theta(t+kP)=\theta(t)+2k\pi$ for every integer $k$.
Its increasing $C^1$ coordinate properties and chord reconstruction remain valid.
\end{proposition}
\begin{proof}
\leanok
The chord vector is nonzero because its length is positive and is periodic
because the boundary is periodic. Apply the nonreturn theorem and the
abstract first-return result, then substitute into the integer shift law.
\end{proof}

\begin{definition}[Cyclic angle coordinates]
\label{def:cyclic-angle-coordinates}
\lean{UlamFloating.cyclicAngleCoordinates}
\leanok
\uses{def:sampling}
For a real chord-direction lift $\theta$ and $h=P/n$, set
\[
 x_i(t)=\pi+\theta(t+ih)-\theta(t+(i+1)h),\qquad i\in\Z/n\Z.
\]
The coordinate at $i$ corresponds to the shared vertex of edges $i$ and
$i+1$. Geometric range conditions and differential equations are proved
separately; they are not part of the definition.
\end{definition}

\begin{lemma}[Regularity, periodicity, and cyclic shift of angle coordinates]
\label{lem:cyclic-angle-api}
\lean{UlamFloating.periodic_angle_difference,UlamFloating.cyclicAngleCoordinates_contDiff,UlamFloating.cyclicAngleCoordinates_periodic,UlamFloating.cyclicAngleCoordinates_rotate}
\leanok
\uses{def:cyclic-angle-coordinates,lem:finite-cycle}
If $\theta$ is $C^1$ and $\theta(t+P)=\theta(t)+A$, each $x_i$ is $C^1$
and $P$-periodic, and $x_{i+1}(t)=x_i(t+P/n)$, including cyclic wraparound.
\end{lemma}
\begin{proof}
\leanok
The fixed increment $A$ cancels in a difference of shifted angles.
Compose with parameter translations and apply finite cyclic sampling.
\end{proof}

\begin{proposition}[Exact cyclic angle sum]
\label{prop:cyclic-angle-sum}
\lean{UlamFloating.cyclicAngleCoordinates_sum}
\leanok
\uses{def:cyclic-angle-coordinates}
If $n>0$ and $\theta(t+P)=\theta(t)+A$, then $\sum_i x_i(t)=n\pi-A$.
For a one-turn chord lift this is $(n-2)\pi$, hence $3\pi$ for five samples.
\end{proposition}
\begin{proof}
\leanok
The sum of angle differences telescopes from $\theta(t)$ to
$\theta(t+P)$. The remaining $n$ constant summands contribute $n\pi$.
\end{proof}

\begin{proposition}[Convex range of the real angle coordinates]
\label{prop:chord-angle-range}
\lean{UlamFloating.IsUnitSpeedBoundary.chord_angle_increment_lt_pi,UlamFloating.IsUnitSpeedBoundary.angle_coordinate_pos_lt_pi}
\leanok
\uses{prop:global-chord-angle,lem:nonparallel-chords}
For $0<h$ and $2h<P$, an increasing chord-direction lift satisfies
$0<\theta(t+h)-\theta(t)<\pi$. Thus
$0<\pi+\theta(t)-\theta(t+h)<\pi$.
\end{proposition}
\begin{proof}
\leanok
Strict monotonicity gives the lower bound. If the direction advanced by
at least $\pi$, the inverse coordinate at $\theta(t)+\pi$ would produce
an opposite chord vector before or at $t+h$. The chords would then be
interlaced or adjacent and parallel, contradicting convexity.
\end{proof}

\begin{proposition}[Geometric extraction of periodic cyclic angle coordinates]
\label{prop:cyclic-angle-extraction}
\lean{UlamFloating.IsUnitSpeedBoundary.exists_cyclic_angle_coordinates}
\leanok
\uses{prop:chord-one-turn,lem:cyclic-angle-api,prop:cyclic-angle-sum,prop:chord-angle-range}
For every $n\geq3$, constant-chord boundary data at increment $P/n$ admit
a global chord-direction lift and cyclic coordinates $x_i$ that are $C^1$,
$P$-periodic, lie in $(0,\pi)$, sum to $(n-2)\pi$, and satisfy
$x_{i+1}(t)=x_i(t+P/n)$. The lift reconstructs the original chord vectors.
\end{proposition}
\begin{proof}
\leanok
Since $0<P/n$ and $2P/n<P$, construct the one-turn lift. Apply the
regularity, periodicity, angle-sum, range, and cyclic-shift lemmas.
The pentagon differential equation is established separately below.
\end{proof}

\begin{lemma}[Derivative of a planar rotation]
\label{lem:rotation-derivative}
\lean{UlamFloating.hasDerivAt_rotation,UlamFloating.angular_derivative_eq}
\leanok
\uses{lem:rotation-reconstruction}
For any differentiable real angle, $(R_{\theta(t)}v)'=
\theta'(t)JR_{\theta(t)}v$. If a nonzero curve satisfies both
$e'=aJe$ and $e=R_\theta e(0)$, then $\theta'=a$.
These statements hold in any oriented two-dimensional real inner product space.
\end{lemma}
\begin{proof}
\leanok
Differentiate the sine and cosine coefficients. Uniqueness of the derivative
and injectivity of scalar multiplication by the nonzero vector $Je$ identify
the angular coefficient.
\end{proof}

\begin{proposition}[Exact angular speed of a boundary chord lift]
\label{prop:chord-angle-speed}
\lean{UlamFloating.IsUnitSpeedBoundary.hasDerivAt_chord_rotation,UlamFloating.IsUnitSpeedBoundary.chord_angle_deriv,UlamFloating.IsUnitSpeedBoundary.chord_area_pos_of_angle_deriv_pos}
\leanok
\uses{lem:second-carousel,lem:rotation-derivative,prop:kinematic-extraction}
In either orientation, a constant chord satisfies
$e'=2\omega(\gamma',e)Je/\ell^2$. Every differentiable real lift
reconstructing these chords therefore satisfies
$\theta'=2\omega(\gamma',e)/\ell^2$.
If $\theta'>0$, then $\omega(\gamma',e)>0$ in that same orientation.
\end{proposition}
\begin{proof}
\leanok
Differentiate the endpoints and apply the two local carousel laws.
Compare this derivative with the derivative of the reconstructed rotation.
The positive chord length allows division, and its square is positive.
\end{proof}

\begin{lemma}[Positive unit-vector and tangent--chord angles]
\label{lem:positive-chord-angle}
\lean{UlamFloating.exists_positive_unit_angle,UlamFloating.reflection_eq_rotation,UlamFloating.exists_positive_chord_angle}
\leanok
\uses{lem:second-carousel}
For unit vectors $u,w$ with $\omega(u,w)>0$, there is
$\alpha\in(0,\pi)$ such that
$\cos\alpha=\ip uw$, $\sin\alpha=\omega(u,w)$, and
$w=R_\alpha u$. Reflection of $u$ in the line through $w$ is $R_\alpha w$.
Consequently the carousel laws with $\norm e=\ell>0$ and $\norm u=1$
give $e/\ell=R_\alpha u$, $v=R_\alpha(e/\ell)$, and
$\sin\alpha=\omega(u,e)/\ell$ whenever $\omega(u,e)>0$.
\end{lemma}
\begin{proof}
\leanok
Use $\alpha=\arccos\ip uw$. The identity
$\ip uw^2+\omega(u,w)^2=1$ and positive area determine the sine and
put the angle strictly in $(0,\pi)$. The oriented orthogonal decomposition
reconstructs $w$. Substitute in the reflection formula, then normalize
a positive-length chord to a unit vector.
\end{proof}

\begin{lemma}[Real addition of consecutive tangent--chord angles]
\label{lem:tangent-angle-add}
\lean{UlamFloating.rotation_angle_unique,UlamFloating.tangent_chord_angle_add}
\leanok
\uses{lem:positive-chord-angle,lem:rotation-integer-turns}
Two angles in $(0,2\pi)$ inducing the same rotation of a nonzero vector
are equal as real numbers. In particular, if consecutive tangent--chord
angles $\alpha,\beta$ lie in $(0,\pi)$ and the chord direction increment
$\delta$ lies in $(0,\pi)$, the reflection identities imply
$\alpha+\beta=\delta$.
\end{lemma}
\begin{proof}
\leanok
Compose one rotation with the inverse of the other. The difference is
an integer multiple of $2\pi$ and lies strictly between $-2\pi$ and $2\pi$,
so it vanishes. Compose the two tangent--chord rotations and apply this
uniqueness result; no full-turn ambiguity remains.
\end{proof}

\begin{proposition}[Periodic tangent angles and the geometric sine law]
\label{prop:boundary-tangent-angles}
\lean{UlamFloating.IsUnitSpeedBoundary.exists_tangent_chord_angles,UlamFloating.IsUnitSpeedBoundary.exists_tangent_chord_angles_with_rotation}
\leanok
\uses{prop:chord-angle-speed,lem:positive-chord-angle,lem:tangent-angle-add,prop:chord-angle-range}
Suppose $0<h<P/2$ and an increasing $C^1$ chord lift $\theta$ has
positive derivative. There is a $P$-periodic function
$\alpha:\mathbb R\to(0,\pi)$ such that
\[
 \alpha(t)+\alpha(t+h)=\theta(t+h)-\theta(t),\qquad
 \theta'(t)=\frac2\ell\sin\alpha(t).
\]
Both identities follow from the original boundary hypotheses. The stronger
API also retains the geometric identity
$e(t)/\ell=R_{\alpha(t)}\gamma'(t)$.
\end{proposition}
\begin{proof}
\leanok
Construct the positive tangent--chord angle at each time. Its cosine is
the inner product of the unit tangent and normalized chord. These vectors
are periodic; injectivity of cosine on $[0,\pi]$ proves periodicity of
$\alpha$. The reflection law and the convex increment bound give the
consecutive-sum identity. Substitute the sine formula in the exact
angular speed.
\end{proof}

\begin{lemma}[Solving the five tangent-angle equations]
\label{lem:five-angle-relations}
\lean{UlamFloating.five_tangent_angle_eq,UlamFloating.five_chord_angle_speed,UlamFloating.hasDerivAt_five_angle_coordinate}
\leanok
\uses{prop:boundary-tangent-angles,prop:chord-one-turn}
Let $h=P/5$ and $x(t)=\pi+\theta(t)-\theta(t+h)$.
The periodic tangent-angle system and the one-turn identity imply
\[
 \alpha(t)=\pi-\bigl(\theta(t+2h)-\theta(t+h)\bigr)
                 -\bigl(\theta(t+4h)-\theta(t+3h)\bigr).
\]
For a sine speed law $\theta'=k\sin\alpha$, it follows that
\[
 x'(t)=k\bigl(\sin(x(t+2h)+x(t+4h))
                  -\sin(x(t+h)+x(t+3h))\bigr).
\]
The Lean lemmas allow any real scaling coefficient $k$.
\end{lemma}
\begin{proof}
\leanok
Solve the five real linear equations, using $\alpha(t+5h)=\alpha(t)$
and $\theta(t+5h)=\theta(t)+2\pi$. Rewrite the resulting angle as
$x(t+h)+x(t+3h)-\pi$. Apply $\sin(y-\pi)=-\sin y$ and differentiate
the defining consecutive difference for $x$.
\end{proof}

\begin{definition}[The Euclidean curve of sampled angles]
\label{def:cyclic-angle-curve}
\lean{UlamFloating.cyclicAngleCurve}
\leanok
\uses{def:cyclic-angle-coordinates}
Bundle the cyclic angle coordinates into the Euclidean vector
$X(t)=(x_i(t))_{i\in\mathbb Z/n\mathbb Z}$, with its Euclidean norm.
\end{definition}

\begin{lemma}[Regularity, periodicity, and arbitrary cyclic shifts]
\label{lem:angle-curve-api}
\lean{UlamFloating.boundarySamples_add_nat,UlamFloating.cyclicAngleCurve_contDiff,UlamFloating.cyclicAngleCurve_periodic,UlamFloating.cyclicAngleCoordinates_add_nat}
\leanok
\uses{def:cyclic-angle-curve,lem:cyclic-angle-api,lem:finite-cycle}
The curve $X$ is $C^1$ and $P$-periodic. For every natural number $k$,
$x_{i+k}(t)=x_i(t+kP/n)$, including all cyclic wraps.
The underlying sampling-shift lemma works for any periodic function into
any target type and every nonempty finite cycle.
\end{lemma}
\begin{proof}
\leanok
Apply coordinatewise regularity and periodicity. Induct on $k$ using the
one-step sample rotation identity; specialize to the periodic angle difference.
\end{proof}

\begin{theorem}[The geometric pentagon differential equation]
\label{thm:geometric-pentagon-ode}
\lean{UlamFloating.IsUnitSpeedBoundary.pentagon_angle_deriv}
\leanok
\uses{lem:five-angle-relations,lem:angle-curve-api,prop:boundary-tangent-angles}
For the actual five sampled coordinates of a once-turning positive chord
lift of a constant one-fifth chord boundary,
\[
 X'(t)=\frac2\ell F(X(t)),\qquad
 F(X)_i=\sin(X_{i+2}+X_{i+4})-\sin(X_{i+1}+X_{i+3}).
\]
This equation is a conclusion of the boundary geometry.
\end{theorem}
\begin{proof}
\leanok
Construct the periodic tangent angles and apply the scalar five-angle
derivative formula at each sampled time. Arbitrary cyclic shift identities
convert time translations to finite-index additions. Coordinate projections
and uniqueness of derivatives identify the Euclidean derivative.
\end{proof}

\begin{proposition}[Full geometric extraction of the pentagon angle solution]
\label{prop:pentagon-angle-solution}
\lean{UlamFloating.IsUnitSpeedBoundary.exists_pentagon_angle_solution}
\leanok
\uses{thm:geometric-pentagon-ode,prop:chord-one-turn,prop:cyclic-angle-sum,prop:chord-angle-range,lem:angle-curve-api}
A constant one-fifth chord boundary admits an increasing real chord lift
reconstructing its chord vectors, whose sampled curve $X$ is $C^1$,
$P$-periodic, has coordinates in $(0,\pi)$ summing to $3\pi$, satisfies
$X_{i+1}(t)=X_i(t+P/5)$, and solves $X'=(2/\ell)F(X)$.
The formal statement retains the lift and its reconstruction of the given
boundary, rather than merely asserting existence of an unrelated solution.
\end{proposition}
\begin{proof}
\leanok
Choose the one-turn lift and apply the proved regularity, range, sum,
shift, and geometric differential-equation results to that same lift.
\end{proof}

\begin{lemma}[Constant sampled angles determine the increment]
\label{lem:constant-sampled-angles}
\lean{UlamFloating.cyclicAngleCoordinates_eq_initial_of_constant,UlamFloating.angle_increment_eq_of_constant_cyclic_curve,UlamFloating.cyclicAngleCoordinates_value_of_constant}
\leanok
\uses{def:cyclic-angle-curve,prop:cyclic-angle-sum}
For any nonempty sampling cycle, if $X(t)$ is constant in time and
$\theta(t+P)=\theta(t)+A$, then every sampled coordinate is
$\pi-A/n$ and $\theta(t+P/n)-\theta(t)=A/n$.
No geometric assumptions on the underlying functions are needed.
\end{lemma}
\begin{proof}
\leanok
Every sample is a time translate of sample zero. Time constancy therefore
makes all coordinates equal. Their telescoping sum determines the common
value and the chord increment.
\end{proof}

\begin{proposition}[Uniformly rotating velocity gives a circle]
\label{prop:rotating-velocity-circle}
\lean{UlamFloating.rotating_velocity_center_eq,UlamFloating.exists_circle_of_rotating_unit_velocity,UlamFloating.rotating_velocity_position_eq,UlamFloating.IsUnitSpeedBoundary.isDisk_of_rotating_velocity}
\leanok
\uses{lem:rotation-reconstruction,prop:disk}
In an oriented two-dimensional real inner product space, suppose
$\gamma'=v$ and $v'=aJv$ with constant $a\ne0$. Then
$c=\gamma(t)+a^{-1}Jv(t)$ is independent of $t$, and
\[
 \gamma(t)-c=R_{at}\bigl(\gamma(0)-c\bigr).
\]
If $\norm{v(t)}=1$ and $a>0$, the curve lies on the circle of radius
$a^{-1}$ about $c$. A body with such a unit-speed boundary is a disk.
\end{proposition}
\begin{proof}
\leanok
Differentiate the proposed center and use $J^2=-I$. Its derivative vanishes.
Reconstruct $v$ from the rotation equation and substitute in the center
identity to recover the position. Norm preservation gives the radius;
the frontier-to-ball theorem identifies the body.
\end{proof}

\begin{lemma}[Rotation distances and regular polygons]
\label{lem:rotation-distance-regularity}
\lean{UlamFloating.rotation_dist_sq,UlamFloating.isRegularPolygon_of_rotation}
\leanok
\uses{def:polygon-geometry}
For any vector $v$ in an oriented Euclidean plane,
\[
 \norm{R_a v-R_b v}^2=2\norm v^2\bigl(1-\cos(a-b)\bigr).
\]
Thus $p_i=c+R_{\phi+2\pi i/n}v$, with $n\geq3$ and $v\ne0$, is a
regular polygon in the full distance-table sense used by this blueprint.
\end{lemma}
\begin{proof}
\leanok
Expand the squared norm using the inner product. The sine and cosine
addition formulas give the distance identity; norm preservation gives
the common positive circumradius.
\end{proof}

\begin{proposition}[One turn of a uniformly rotating boundary]
\label{prop:uniform-rotation-one-turn}
\lean{UlamFloating.IsUnitSpeedBoundary.rotation_period_eq_two_pi}
\leanok
\uses{lem:angle-first-return,lem:boundary-embedding}
If a simple periodic boundary satisfies
$\gamma(t)-c=R_{at}v$ with $a>0$ and $v\ne0$, then $aP=2\pi$.
\end{proposition}
\begin{proof}
\leanok
Apply the first-return theorem to the increasing angle $t\mapsto at$.
Periodicity gives return at $P$. Injectivity on one half-open traversal
excludes an earlier return, forcing exactly one full turn.
\end{proof}

\begin{proposition}[Regular samples of a uniformly rotating boundary]
\label{prop:uniform-rotation-regular-samples}
\lean{UlamFloating.IsUnitSpeedBoundary.boundarySamples_regular_of_rotating_velocity}
\leanok
\uses{prop:rotating-velocity-circle,prop:uniform-rotation-one-turn,lem:rotation-distance-regularity}
For every $n\geq3$, rational samples of a unit-speed boundary whose tangent
rotates at a constant positive rate form a regular polygon in their original
cyclic ordering.
\end{proposition}
\begin{proof}
\leanok
Reconstruct the position about its fixed center and use $aP=2\pi$.
At sample $i$ the rotation angle is $at+2\pi i/n$.
Apply the complete rotation distance formula, which excludes star ordering.
\end{proof}

\begin{proposition}[A constant pentagon angle curve gives uniform tangent rotation]
\label{prop:constant-fifth-velocity}
\lean{UlamFloating.IsUnitSpeedBoundary.fifth_velocity_rotation_of_constant_angles}
\leanok
\uses{lem:constant-sampled-angles,prop:boundary-tangent-angles,lem:five-angle-relations}
If the actual sampled-angle curve of a one-fifth chord boundary is constant,
then its original unit tangent satisfies
\[
 \gamma''(t)=aJ\gamma'(t),\qquad
 a=\frac2\ell\sin(\pi/5)>0.
\]
Existence of the second derivative is a conclusion, despite the initial
boundary hypotheses requiring only $C^1$ regularity.
\end{proposition}
\begin{proof}
\leanok
The constant sampled vector forces all chord increments to be $2\pi/5$.
The five tangent-angle equations then give $\alpha=\pi/5$ and the sine
speed law gives $\theta'=a$. The retained tangent--chord rotation identity
expresses $\gamma'$ as a fixed vector rotated by $\theta-\pi/5$.
Differentiate this expression to obtain the stated equation.
\end{proof}

\begin{proposition}[The constant pentagon case is rigid]
\label{prop:constant-fifth-rigidity}
\lean{UlamFloating.IsUnitSpeedBoundary.isDisk_of_constant_fifth_angles,UlamFloating.IsUnitSpeedBoundary.regular_pentagons_of_constant_angles}
\leanok
\uses{prop:constant-fifth-velocity,prop:rotating-velocity-circle,prop:uniform-rotation-regular-samples}
A constant geometric five-angle curve forces both disk rigidity and
regularity of every sampled pentagon in its specified cyclic order.
The supplied angle lift is the positive once-turning lift of the boundary
chords; it is not an arbitrary solution of the abstract angle equation.
\end{proposition}
\begin{proof}
\leanok
The tangent rotates at the positive constant rate
$(2/\ell)\sin(\pi/5)$. Apply the circle reconstruction and regular-sampling
theorems. This proves the constant branch of the pentagon obstruction.
\end{proof}

\section{Invariants and the disk bridge}

\begin{lemma}[Bilinear area-rate cancellation]
\label{lem:area-rate}
\lean{UlamFloating.sum_area_rate_eq_zero}
\leanok
For a skew real bilinear form $B$, a permutation $s$ of a finite index
set, and vectors with $v_i+v_{s(i)}=a_i(p_{s(i)}-p_i)$,
\[
 \sum_i\bigl(B(v_i,p_{s(i)})+B(p_i,v_{s(i)})\bigr)=0.
\]
\end{lemma}
\begin{proof}
\leanok
Skewness gives $B(x,x)=0$. Pair each velocity sum with its edge.
The summand becomes $B(v_i,p_i)-B(v_{s(i)},p_{s(i)})$.
Reindex the second sum by $s$. This needs neither convexity nor equal
edge lengths, and works in any dimension.
\end{proof}

\begin{definition}[Signed polygon moments and area centroid]
\label{def:moments}
\lean{UlamFloating.polygonSignedArea,UlamFloating.polygonFirstMoment,UlamFloating.polygonCentroid,UlamFloating.planarArea,UlamFloating.laminaCentroid}
\leanok
\uses{def:polygon-geometry}
For an orientation and a finite successor permutation, define signed area
$A$ and signed first moment $M$ by the formulas below, and signed centroid
as $A^{-1}M$. For a measurable planar set, define its area using Lebesgue
measure and its lamina centroid as the normalized integral of position.
Geometric applications establish positive finite area or nonzero signed
area before using these normalized expressions.
\end{definition}

\begin{proposition}[Signed-area conservation]
\label{prop:area}
\lean{UlamFloating.sum_bilinear_apply_const_of_midpoint_parallel,UlamFloating.IsCarousel.polygonSignedArea_eq,UlamFloating.carousel_area}
\leanok
\uses{def:carousel,lem:area-rate,def:moments}
For differentiable moving vertices in a real normed space, the cyclic sum
of a continuous skew bilinear form is constant under the midpoint-parallel
law. In particular, the signed area of a planar carousel is constant.
\end{proposition}
\begin{proof}
\leanok
Differentiate each bilinear summand and sum the derivatives. The area-rate
identity makes the derivative zero everywhere. The mean value theorem
gives constancy. Convexity and equal edge lengths are not needed here.
\end{proof}

\begin{lemma}[First-moment edge identities]
\label{lem:moment-rate}
\lean{UlamFloating.areaForm_smul_eq_rotation,UlamFloating.areaForm_smul_eq_velocity_difference,UlamFloating.moment_edge_rate,UlamFloating.sum_moment_rate_eq_zero}
\leanok
\uses{lem:orthogonal,def:polygon-geometry}
Write $\omega$ for the oriented area form and $J$ for the positive
right-angle rotation. For an edge $e=q-p$ and velocities $u,v$ satisfying
$u+v=ae$ and $\ip e{v-u}=0$, the derivative-shaped expression is
\[
 (\omega(u,q)+\omega(p,v))(p+q)+\omega(p,q)(u+v)
 =2\omega(u,p)p-2\omega(v,q)q
   +\tfrac12\norm e^2(Jv-Ju).
\]
For equal-length edges its sum over any finite successor permutation is zero.
\end{lemma}
\begin{proof}
\leanok
Use the planar identity
$\omega(u,e)e=\norm e^2 Ju-\ip eu\,Je$.
The two velocity laws give $2\ip eu=a\norm e^2$, hence
$\omega(u,e)e=-\tfrac12\norm e^2(Jv-Ju)$.
Bilinearity and the midpoint law give the displayed edge identity.
For common edge length $\ell$, each summand is $F_i-F_{s(i)}$, where
$F_i=2\omega(v_i,p_i)p_i-\tfrac12\ell^2Jv_i$; reindex to cancel.
\end{proof}

\begin{proposition}[First moment and signed centroid]
\label{prop:invariants}
\lean{UlamFloating.IsCarousel.sum_areaForm_smul_eq,UlamFloating.IsCarousel.polygonFirstMoment_eq,UlamFloating.IsCarousel.polygonCentroid_eq,UlamFloating.carousel_firstMoment,UlamFloating.carousel_centroid}
\leanok
\uses{def:carousel,prop:area,lem:moment-rate,def:moments}
The signed first moment of a planar carousel is constant.
Its signed centroid is fixed. These statements quantify over any finite
successor permutation; the moment proof works in every oriented real
inner product space of dimension two. Geometric interpretation of the
centroid requires nonzero signed area.
\end{proposition}
\begin{proof}
\leanok
Differentiate finite sums using mathlib's oriented area form:
\[
 A=\tfrac12\sum_i\det(p_i,p_{s(i)}),\qquad
 M=\tfrac16\sum_i\det(p_i,p_{s(i)})(p_i+p_{s(i)}).
\]
The first-moment edge identity makes $M'=0$. The mean value theorem gives
constancy, and constancy of $A$ and $M$ gives constancy of $A^{-1}M$.
\end{proof}

\begin{lemma}[Signed area as a sum of edge triangles]
\label{lem:edge-origin}
\lean{UlamFloating.sum_bilinear_edge_origin}
\leanok
For every skew real bilinear form $B$, finite successor permutation $s$,
vertices $p_i$, and origin $c$,
\[
 \sum_i B(p_{s(i)}-p_i,c-p_i)=\sum_i B(p_i,p_{s(i)}).
\]
This identity holds in any real vector space.
\end{lemma}
\begin{proof}
\leanok
Expand each summand by bilinearity and skew symmetry. The terms involving
$c$ cancel after reindexing the finite sum by the successor permutation.
\end{proof}

\begin{proposition}[Nonzero signed area of a strictly convex polygon]
\label{prop:polygon-area-nonzero}
\lean{UlamFloating.polygonSignedArea_neg_orientation,UlamFloating.polygonSignedArea_pos_of_edge_sides,UlamFloating.IsStrictlyConvexPolygon.polygonSignedArea_ne_zero}
\leanok
\uses{lem:edge-origin,def:polygon-geometry,def:moments}
The signed area is positive in an orientation satisfying the strict
edge-side inequalities. Reversing orientation negates it. Thus a strictly
convex polygon has nonzero signed area in either orientation.
\end{proposition}
\begin{proof}
\leanok
Choose a vertex as the origin in the edge-triangle identity. Every
edge contributes nonnegatively, and an edge not incident to that vertex
contributes positively. Such an edge exists because $n\geq3$.
There are exactly two orientations, related by negation of the area form.
\end{proof}

\begin{lemma}[Alternating area-vector identity]
\label{lem:area-vector-triple}
\lean{UlamFloating.areaForm_smul_three}
\leanok
In every oriented Euclidean plane,
\[
 \omega(a,b)c=\omega(a,c)b+\omega(c,b)a.
\]
\end{lemma}
\begin{proof}
\leanok
For nonzero $a$, expand $b$ and $c$ in the basis $(a,Ja)$ and use
bilinearity. If $a=0$, all terms vanish.
\end{proof}

\begin{proposition}[Translation of polygon moments]
\label{prop:moment-translation}
\lean{UlamFloating.sum_areaForm_translate,UlamFloating.sum_moment_translate,UlamFloating.polygonSignedArea_translate,UlamFloating.polygonFirstMoment_translate,UlamFloating.polygonCentroid_translate}
\leanok
\uses{lem:area-vector-triple,def:moments}
For any finite successor permutation, translating every vertex by $c$
preserves signed area $A$ and changes signed first moment $M$ to $M+Ac$.
If $A\ne0$, the signed centroid translates by $c$.
\end{proposition}
\begin{proof}
\leanok
Expand the area form and cancel differences under the successor
permutation. For first moments, the alternating area-vector identity
reduces the extra terms to three times the shoelace sum times $c$,
plus a telescoping sum. Multiply by $1/6$ and divide by $A$.
\end{proof}

\begin{lemma}[Centered moment of an equilateral cyclic polygon]
\label{lem:cyclic-moment-zero}
\lean{UlamFloating.areaForm_smul_add_of_norm_sq_eq,UlamFloating.sum_moment_eq_zero_of_norm_sq_eq}
\leanok
\uses{lem:moment-rate}
If $\norm a^2=\norm b^2=R$ and $\norm{b-a}=\ell$, then
\[
 \omega(a,b)(a+b)=(2R-\ell^2/2)(Ja-Jb).
\]
Thus equal-radius vertices with equal-length edges have zero first moment
about the origin, for any finite successor permutation.
\end{lemma}
\begin{proof}
\leanok
The squared edge length gives $\ip ab=R-\ell^2/2$. Apply the planar
area-rotation identity to each endpoint. Sum the resulting differences
of rotated vertices; the common coefficient factors out and the sum vanishes.
\end{proof}

\begin{proposition}[Circumcenter as signed centroid]
\label{prop:cyclic-centroid}
\lean{UlamFloating.polygonCentroid_eq_of_equidistant}
\leanok
\uses{prop:moment-translation,lem:cyclic-moment-zero}
An equilateral cyclic polygon with nonzero signed area has signed
centroid equal to its circumcenter. Convexity and regular labeling are
unnecessary for this statement.
\end{proposition}
\begin{proof}
\leanok
Translate the circumcenter to the origin. The centered first moment is
zero. Translate the normalized centroid back using nonzero signed area.
\end{proof}

\begin{proposition}[Geometric centroid identification: statement formalized]
\label{prop:geometric-centroid}
\lean{UlamFloating.polygonCentroidStatement}
\uses{def:moments,def:polygon-geometry,prop:polygon-area-nonzero}
A strictly convex polygon has nonzero signed area, and its signed centroid
equals the lamina centroid of the convex hull of its vertices, in either
orientation.
\end{proposition}
The lamina centroid is generally not the average of the vertices.
Nonvanishing of signed area is proved above. The centroid identification
still requires a triangulation and integral proof.
Boundary data permit both orientations, so positivity of signed area
must not be assumed silently.

\begin{lemma}[Reconstruction from a spherical frontier]
\label{lem:ball-frontier}
\lean{UlamFloating.subset_closedBall_of_frontier_subset,UlamFloating.eq_closedBall_of_frontier_subset_sphere}
\leanok
In a nontrivial real normed space, a nonempty compact set lies in any
closed ball containing its frontier. If it has nonempty interior and its
frontier is contained in a sphere of positive radius, it equals the closed
ball. Convexity is not required.
\end{lemma}
\begin{proof}
\leanok
Maximize distance from the ball center on the compact set. A maximizer
cannot be an interior point: the set lies in the closed ball with this
maximal radius, whose interior is the open ball. Thus the maximum occurs
on the frontier, proving containment.
For the second claim, the open ball contains no frontier point. Partition
it into its intersections with the interior of the set and the complement
of the set. These are disjoint open pieces, and the first is nonempty.
Connectedness forces the entire open ball into the set. Taking closures
gives the opposite inclusion.
\end{proof}

\begin{proposition}[Circle boundary to disk]
\label{prop:disk}
\lean{UlamFloating.IsUnitSpeedBoundary.eq_closedBall,UlamFloating.circle_boundary_disk,UlamFloating.IsUnitSpeedBoundary.isDisk_of_norm_sq}
\leanok
\uses{def:boundary,def:perimetral,lem:ball-frontier}
If boundary data satisfy $\norm{\gamma(t)-c}=r>0$ for all $t$, then
$K=\overline B(c,r)$. Equivalently, a positive constant squared distance
$\norm{\gamma(t)-c}^2=q>0$ gives a disk of radius $\sqrt q$.
\end{proposition}
\begin{proof}
\leanok
The range condition turns the norm equality into containment of the
frontier in the sphere. Apply the compact-set reconstruction theorem to
obtain equality of bodies.
\end{proof}

\section{A reusable period obstruction}

\begin{lemma}[Fourier energy on one interval]
\label{lem:fourier-energy}
\lean{UlamFloating.summable_sq_fourierCoeff,UlamFloating.tsum_sq_fourierCoeff_continuous,UlamFloating.hasSum_sq_fourierCoeffOn}
\leanok
For a continuous complex-valued function on $[0,T]$, $T>0$, with
matching endpoint values, the squared Fourier coefficients are summable and
\[
 \sum_{k\in\Z}|\widehat f(k)|^2=\frac1T\int_0^T|f(t)|^2\,dt.
\]
The normalization agrees with mathlib's \texttt{fourierCoeffOn}.
\end{lemma}
\begin{proof}
\leanok
Lift the function continuously to the additive circle. Apply mathlib's
Fourier Hilbert basis and Parseval identity, then convert its normalized
Haar integral to an interval integral. The continuous-function version
also removes the intermediate almost-everywhere quotient.
\end{proof}

\begin{lemma}[Nonzero Fourier-mode energy estimate]
\label{lem:fourier-mode}
\lean{UlamFloating.fourierCoeffOn_energy_le}
\leanok
If $f'=g$ on $[0,T]$, $g$ is interval integrable, and $f(0)=f(T)$,
then for every nonzero integer $k$,
\[
 (2\pi)^2|\widehat f(k)|^2\leq T^2|\widehat g(k)|^2.
\]
\end{lemma}
\begin{proof}
\leanok
Integration by parts gives $2\pi i k\widehat f(k)=T\widehat g(k)$.
Take squared norms and use $k^2\geq1$.
\end{proof}

\begin{proposition}[Complex periodic Wirtinger inequality]
\label{prop:wirtinger-complex}
\lean{UlamFloating.wirtinger_complex}
\leanok
\uses{lem:fourier-energy,lem:fourier-mode}
Suppose $f,g:[0,T]\to\mathbb C$ are continuous, $f'=g$, both have
matching endpoint values, and $\int_0^T f=0$. Then
\[
 (2\pi)^2\int_0^T|f|^2\leq T^2\int_0^T|g|^2.
\]
\end{proposition}
\begin{proof}
\leanok
The mean-zero hypothesis removes the zero Fourier coefficient. Sum the
nonzero-mode estimates using summability and Parseval, then cancel the
positive normalization factor $1/T$.
\end{proof}

\begin{proposition}[Finite-dimensional periodic Wirtinger inequality]
\label{prop:wirtinger-euclidean}
\lean{UlamFloating.wirtinger_euclidean}
\leanok
\uses{prop:wirtinger-complex}
The same inequality, with squared Euclidean norms, holds for functions
into any finite-dimensional real Euclidean space under the same
continuity, derivative, endpoint, and mean-zero hypotheses.
\end{proposition}
\begin{proof}
\leanok
Compose with each coordinate projection and the real embedding into
$\mathbb C$. Apply the scalar theorem and sum over the finite coordinate
set. Finite sums commute with the interval integral, and the sum of
squared coordinate norms is the squared Euclidean norm.
\end{proof}

\begin{lemma}[Periodicity, mean, and positive velocity energy]
\label{lem:periodic-energy}
\lean{UlamFloating.periodic_deriv,UlamFloating.integral_deriv_eq_zero_of_periodic,UlamFloating.integral_norm_sq_pos_of_periodic,UlamFloating.integral_deriv_norm_sq_pos}
\leanok
Derivatives preserve periodicity. The derivative of a periodic $C^1$
curve has zero integral over a period. A nonzero continuous periodic
function has positive squared-norm integral over a positive period;
in particular, so does the velocity of a nonconstant $C^1$ curve.
These lemmas apply to general real normed spaces, with completeness
required for the vector-valued fundamental theorem of calculus.
\end{lemma}
\begin{proof}
\leanok
Differentiate the time-shift identity for periodicity and apply the
fundamental theorem of calculus for the mean. A nonzero value can be
moved into one period by periodicity; continuity makes its energy
integral positive. An identically zero derivative would make the
original curve constant by the mean value theorem.
\end{proof}

\begin{proposition}[Period bound from a derivative estimate]
\label{prop:period-bound}
\lean{UlamFloating.period_ge_of_derivative_bound,UlamFloating.period_bound}
\leanok
\uses{prop:wirtinger-euclidean,lem:periodic-energy}
Let $x:\R\to\R^d$ be nonconstant, $C^2$, and $T$-periodic, with $T>0$.
If $L>0$ and $\norm{x''(t)}\leq L\norm{x'(t)}$ for all $t$, then
$T\geq 2\pi/L$.
\end{proposition}
\begin{proof}
\leanok
Put $y=x'$. Its integral over a period is zero, and it is periodic.
The vector-valued periodic Wirtinger inequality gives
\[
 0<\int_0^T\norm y^2\leq
 \frac{T^2}{4\pi^2}\int_0^T\norm{y'}^2\leq
 \frac{T^2L^2}{4\pi^2}\int_0^T\norm y^2.
\]
Positivity follows from nonconstancy and continuity of the derivative.
Cancel the positive integral and take nonnegative square roots.
The stronger library theorem only requires the derivative estimate on
$[0,T]$ and allows any finite coordinate index type.
\end{proof}

\begin{proposition}[Period obstruction for an autonomous equation]
\label{prop:ode-period-bound}
\lean{UlamFloating.period_ge_of_vectorField_derivative_bound}
\leanok
\uses{prop:period-bound}
Let a nonconstant $C^2$ curve solve $x'=F(x)$ with positive period $T$.
If $F$ is differentiable along the trajectory on $[0,T]$ and
$\norm{DF(x(t))}\leq L$ there, with $L>0$, then $T\geq2\pi/L$.
\end{proposition}
\begin{proof}
\leanok
The chain rule gives $x''=DF(x)x'$. The operator-norm bound yields
$\norm{x''}\leq L\norm{x'}$, so the preceding theorem applies.
\end{proof}

Finite-cycle energy bounds now give the exact pentagon-equation estimate
below. Periodic convex angle coordinates and their five-angle differential
equation are now extracted from boundary data. To complete the fifth-density
proof, the least-period and multiplicity results below must be combined with
a stronger geometric or period obstruction, including exceptional
configurations. Yorke's classical Lipschitz period theorem motivates the
general approach; its more general hypotheses require additional analysis.
\href{https://msekce.karlin.mff.cuni.cz/~prazak/uceni/Odr2/Cviceni/Yorke-period.pdf}{Yorke, 1969, original paper}.

\begin{lemma}[Energy of differences along a permutation]
\label{lem:cycle-difference-energy}
\lean{UlamFloating.sum_sq_sub_perm_le}
\leanok
For a real family on any finite index set and any permutation $s$,
\[
 \sum_i(u_i-u_{s(i)})^2\leq4\sum_i u_i^2.
\]
\end{lemma}
\begin{proof}
\leanok
Use $(a-b)^2\leq2a^2+2b^2$ and reindex one of the sums.
\end{proof}

\begin{lemma}[Five-cycle pair-sum energy]
\label{lem:pentagon-pair-energy}
\lean{UlamFloating.pentagon_pair_energy_identity,UlamFloating.pentagon_pair_energy_le,UlamFloating.pentagon_weighted_difference_energy_le}
\leanok
\uses{lem:cycle-difference-energy}
For five cyclic real coordinates,
\[
 \sum_i(u_{i+2}+u_{i+4})^2+\sum_i(u_i+u_{i+1})^2
 =3\sum_i u_i^2+\Bigl(\sum_i u_i\Bigr)^2.
\]
If $\sum_i u_i=0$, the first energy is at most $3\sum_i u_i^2$.
Multiplying each pair sum by a scalar of absolute value at most one,
then taking cyclic differences, gives energy at most $12\sum_i u_i^2$.
\end{lemma}
\begin{proof}
\leanok
Expand the identity: the two pair families account for all ten unordered
pairs. Drop the nonnegative second energy, apply the diagonal contraction,
and then the permutation-difference estimate.
\end{proof}

\begin{definition}[Five-angle vector field]
\label{def:pentagon-angle-field}
\lean{UlamFloating.pentagonAngleField}
\leanok
For $x\in\R^5$ with cyclic indices, define
\[
 F_i(x)=\sin(x_{i+2}+x_{i+4})-\sin(x_{i+1}+x_{i+3}).
\]
This is the side-length-two angle equation. The following results concern
solutions of this explicitly specified equation. The geometric coordinates
constructed above have not yet been proved to solve it.
\end{definition}

\begin{lemma}[Regularity and zero-sum velocity of the angle field]
\label{lem:pentagon-field-api}
\lean{UlamFloating.pentagonAngleField_apply,UlamFloating.contDiff_pentagonAngleField,UlamFloating.pentagonAngleField_sum}
\leanok
\uses{def:pentagon-angle-field}
The inverse-successor definition agrees with the displayed cyclic formula.
The vector field is smooth (indeed analytic), and $\sum_iF_i(x)=0$ for every $x$.
\end{lemma}
\begin{proof}
\leanok
Check the five cyclic indices. Coordinate sums and sine are analytic.
The two sine sums coincide after reindexing by the inverse successor.
\end{proof}

\begin{lemma}[One step of regularity for an autonomous solution]
\label{lem:ode-regularity}
\lean{UlamFloating.contDiff_two_of_deriv_eq}
\leanok
In any real normed space, a $C^1$ solution of $x'=F(x)$ with $F$ of class
$C^1$ is $C^2$.
\end{lemma}
\begin{proof}
\leanok
The derivative equals the $C^1$ composition $F\circ x$.
Use the characterization of $C^2$ regularity by its first derivative.
\end{proof}

\begin{proposition}[Acceleration estimate for the five-angle equation]
\label{prop:pentagon-acceleration}
\lean{UlamFloating.pentagonAngleField_acceleration_bound}
\leanok
\uses{lem:pentagon-field-api,lem:ode-regularity,lem:pentagon-pair-energy}
Every $C^1$ solution of $x'=F(x)$ satisfies
$\norm{x''}\leq2\sqrt3\norm{x'}$ everywhere.
\end{proposition}
\begin{proof}
\leanok
First obtain $C^2$ regularity. Put $u=x'$ and
$w_i=\cos(x_{i+2}+x_{i+4})(u_{i+2}+u_{i+4})$.
Differentiating the equation gives $x''_i=w_i-w_{i-1}$.
The velocity sum is zero and each cosine has square at most one.
The weighted five-cycle estimate gives $\norm{x''}^2\leq12\norm{x'}^2$.
Take nonnegative square roots.
\end{proof}

\begin{proposition}[Exact lower bound for pentagon-equation periods]
\label{prop:pentagon-period}
\lean{UlamFloating.pentagonAngleField_period_ge,UlamFloating.pentagonAngleField_period_ge_scaled,UlamFloating.pentagonAngleField_period_ge_of_length}
\leanok
\uses{prop:pentagon-acceleration,prop:period-bound}
Every positive period $T$ of a nonconstant $C^1$ solution of $x'=F(x)$
satisfies $T\geq\pi/\sqrt3$. More generally, for $x'=kF(x)$ with $k>0$,
one has $T\geq\pi/(k\sqrt3)$. Thus the side-length-$\ell$ equation
$x'=(2/\ell)F(x)$ has $T\geq\ell\pi/(2\sqrt3)$.
\end{proposition}
\begin{proof}
\leanok
Apply the general period theorem with $L=2\sqrt3$.
For a positive time scale $k$, set $y(s)=x(s/k)$. Then $y'=F(y)$,
$y$ is nonconstant, and $kT$ is a period. This proves the scaled estimate
without an implicit normalization. No decimal bounds or monotonicity
conjectures enter the proof.
\end{proof}

\section{The three rigidity targets}

\begin{lemma}[Triangle carousel center and radius]
\label{lem:triangle-carousel}
\lean{UlamFloating.inner_edges_of_equilateral,UlamFloating.triangle_velocity_sum_eq_zero,UlamFloating.norm_sub_triangle_average_sq,UlamFloating.equilateral_norm_sub_average_sq,UlamFloating.IsCarousel.sum_three_eq,UlamFloating.IsCarousel.triangle_fixed_center}
\leanok
\uses{def:carousel,lem:orthogonal}
Every three-vertex carousel in a real inner product space has a fixed
vertex average $c$, and $\norm{p_i(t)-c}^2=\ell^2/3$.
This needs neither convexity nor a lamina-centroid identification.
\end{lemma}
\begin{proof}
\leanok
Its edge vectors $e,f,g$ sum to zero and have equal norm $\ell$;
their pairwise inner products are $-\ell^2/2$.
Write the midpoint laws as $u+v=ae$, $v+w=bf$, $w+u=cg$.
Pair $v-u=bf-cg$ with $e$, and $w-v=cg-ae$ with $f$.
The rigid-edge derivative and $\ell>0$ imply $a=b=c$, whence
$u+v+w=0$. The mean value theorem fixes the vertex sum.
Finally, in any real inner product space,
\[
 \left\|a-\tfrac13(a+b+c)\right\|^2
 =\frac{2\norm{a-b}^2+2\norm{a-c}^2-\norm{b-c}^2}{9}.
\]
Apply the equal-length identities to obtain the radius.
\end{proof}

\begin{proposition}[Triangle reduction]
\label{prop:triangle}
\lean{UlamFloating.triangle_reduction}
\leanok
\uses{prop:kinematic-extraction,lem:triangle-carousel}
The extracted triangle has a fixed center $c$ and
$\norm{p_i(t)-c}^2=\ell^2/3$.
\end{proposition}
\begin{proof}
\leanok
Apply the derived carousel laws to the three samples and then the
triangle-carousel fixed-center theorem. Neither cyclic polygon convexity
nor a lamina-centroid theorem is needed for this route.
\end{proof}

\begin{lemma}[Assembly for one third: conditional]
\label{lem:third-assembly}
\lean{UlamFloating.IsUnitSpeedBoundary.isDisk_of_triangle_carousel,UlamFloating.triangle_reduction_of_carousel_extraction,UlamFloating.third_rigidity_of_triangle_reduction,UlamFloating.third_rigidity_of_carousel_extraction}
\leanok
\uses{lem:triangle-carousel,prop:disk,def:perimetral}
If the three boundary samples form a carousel, the body is a disk.
Consequently the carousel-extraction statement implies the triangle
reduction and $\mathcal R(1/3)$. The triangle-reduction statement alone
also implies $\mathcal R(1/3)$. These are implications with explicit
premises, not proofs of extraction or unconditional rigidity.
\end{lemma}
\begin{proof}
\leanok
Take the fixed center, evaluate the radius formula at sample zero,
and apply the squared-distance version of the disk bridge.
\end{proof}

\begin{theorem}[Geometric rigidity at density $1/3$]
\label{thm:third}
\lean{UlamFloating.third_rigidity}
\leanok
\uses{prop:triangle,lem:third-assembly}
$\mathcal R(1/3)$ holds.
\end{theorem}
\begin{proof}
\leanok
Use the proved triangle reduction in the final assembly. The positive
squared radius is $\ell^2/3$, and sample zero is the entire boundary curve.
\end{proof}

\begin{lemma}[Equilateral quadrilaterals are parallelograms]
\label{lem:quadrilateral-geometry}
\lean{UlamFloating.inner_sub_midpoint,UlamFloating.norm_sub_midpoint_sq,UlamFloating.eq_or_eq_neg_of_norm_eq_of_smul,UlamFloating.parallelogram_of_equilateral}
\leanok
\uses{lem:midpoint-extraction}
In the plane, an equilateral quadrilateral with distinct opposite
vertices satisfies $a+c=b+d$. Polygon order need not be assumed.
\end{lemma}
\begin{proof}
\leanok
Let $m=(a+c)/2$. Both $b-m$ and $d-m$ are perpendicular to $c-a$,
by the equal-distance identity for a perpendicular bisector. They have
the same norm, by the squared-distance formula for a midpoint.
In dimension two they are parallel, hence equal or opposite. Distinctness
excludes equality, giving $b+d=2m=a+c$.
\end{proof}

\begin{lemma}[Nonparallel parallelogram edges]
\label{lem:square-nondegenerate}
\lean{UlamFloating.exists_smul_of_areaForm_eq_zero,UlamFloating.areaForm_ne_zero_of_norm_eq,UlamFloating.areaForm_ne_zero_of_parallelogram}
\leanok
\uses{lem:planar-decomposition,lem:quadrilateral-geometry}
A vanishing area form with a nonzero first vector implies parallelism.
Equal-length vectors that are neither equal nor opposite have nonzero
oriented area. Thus an equilateral parallelogram with distinct opposite
vertices has nonparallel adjacent edges.
\end{lemma}
\begin{proof}
\leanok
Use planar decomposition to express the second vector as a scalar
multiple of the first. Equal norms force that scalar to be $1$ or $-1$.
For the parallelogram either choice would identify opposite vertices.
\end{proof}

\begin{lemma}[Parallelogram carousel velocity equations]
\label{lem:square-velocity}
\lean{UlamFloating.sum_four_velocity_eq_zero,UlamFloating.opposite_velocities_of_sum_eq_zero,UlamFloating.inner_edges_eq_zero_of_carousel_laws}
\leanok
\uses{lem:square-nondegenerate,lem:orthogonal}
For edges $e,f,-e,-f$ with nonzero oriented area, the four midpoint
equations force total vertex velocity zero. If also the differentiated
parallelogram identity holds, opposite velocities are opposite. The
rigid-edge equations and a nonzero vertex velocity then force $\ip ef=0$.
\end{lemma}
\begin{proof}
\leanok
Write $u_i+u_{i+1}=a_i e_i$. Alternating the four equations gives
$(a_0-a_2)e+(a_3-a_1)f=0$. Pair with the area form to get $a_0=a_2$,
then add the two opposite midpoint equations to get total velocity zero.
Together with $u_0+u_2=u_1+u_3$, this gives
$u_2=-u_0$ and $u_3=-u_1$. Hence
$u_0+u_1=a_0e$ and $u_1-u_0=a_1f$.
The rigid-edge equations give $a_0\ip ef=a_1\ip ef=0$.
If $\ip ef\neq0$, both coefficients vanish and $u_0=0$, a contradiction.
\end{proof}

\begin{proposition}[Four-vertex carousel motion]
\label{prop:square-motion}
\lean{UlamFloating.IsCarousel.four_parallelogram,UlamFloating.IsCarousel.sum_four_deriv_eq_zero,UlamFloating.IsCarousel.sum_four_eq,UlamFloating.IsCarousel.four_edges_inner_eq_zero}
\leanok
\uses{def:carousel,lem:quadrilateral-geometry,lem:square-velocity}
A four-vertex planar carousel whose vertices remain distinct has a fixed
vertex sum and perpendicular adjacent edges at every time.
\end{proposition}
\begin{proof}
\leanok
At every time the equilateral quadrilateral is a nondegenerate parallelogram.
Apply the velocity cancellation, then the mean value theorem to fix the sum.
Differentiate the parallelogram identity and apply the right-angle calculation.
\end{proof}

\begin{lemma}[Square radius and full distance formula]
\label{lem:square-metric}
\lean{UlamFloating.square_norm_sub_average_sq,UlamFloating.square_distance_sq,UlamFloating.square_cosine_table,UlamFloating.regularPolygon_four_of_parallelogram}
\leanok
\uses{def:polygon-geometry}
An equilateral parallelogram with perpendicular adjacent edges has
squared circumradius $\ell^2/2$ about its vertex average. Its squared
pairwise distances are $0$ at equal indices, $2\ell^2$ at opposite
indices, and $\ell^2$ at adjacent indices. It satisfies the full
regular-polygon distance formula in its given cyclic order.
\end{lemma}
\begin{proof}
\leanok
The centered vertices are the four vectors $\pm(e+f)/2$ and $\pm(e-f)/2$.
Use $\norm e=\norm f=\ell$ and $\ip ef=0$ to compute their norms and
all pairwise distances. The sixteen index pairs match
$1-\cos(2\pi(i-j)/4)$ by the exact values of cosine at multiples of
$\pi/2$. Take the positive radius $\sqrt{\ell^2/2}$.
\end{proof}

\begin{proposition}[Square carousel center and regularity]
\label{prop:square-carousel}
\lean{UlamFloating.IsCarousel.square_fixed_center,UlamFloating.IsCarousel.isRegularPolygon_four}
\leanok
\uses{prop:square-motion,lem:square-metric}
A four-vertex planar carousel with distinct vertices consists of regular
squares with one fixed center and squared circumradius $\ell^2/2$.
\end{proposition}
\begin{proof}
\leanok
The fixed vertex average is the center. Apply the radius and complete
distance formulas at each time.
\end{proof}

\begin{proposition}[Square reduction]
\label{prop:square}
\lean{UlamFloating.square_reduction}
\leanok
\uses{prop:kinematic-extraction,lem:boundary-embedding,prop:square-carousel}
Every extracted quadrilateral is a square with a fixed center $c$, and
$\norm{p_i(t)-c}^2=\ell^2/2$.
\end{proposition}
\begin{proof}
\leanok
The boundary assumptions give carousel laws and distinct samples.
Apply the square-carousel theorem, including its regular-polygon conclusion.
Neither cyclic polygon convexity nor a lamina-centroid identification is needed.
\end{proof}

\begin{lemma}[Assembly for one quarter: conditional]
\label{lem:quarter-assembly}
\lean{UlamFloating.quarter_rigidity_of_square_reduction}
\leanok
\uses{prop:disk,def:perimetral}
The square-reduction statement implies $\mathcal R(1/4)$.
The square-reduction statement is an explicit premise of this assembly lemma.
\end{lemma}
\begin{proof}
\leanok
The squared radius is $\ell^2/2>0$. Use sample zero and the disk bridge.
\end{proof}

\begin{theorem}[Geometric rigidity at density $1/4$]
\label{thm:quarter}
\lean{UlamFloating.quarter_rigidity}
\leanok
\uses{prop:square,lem:quarter-assembly}
$\mathcal R(1/4)$ holds.
\end{theorem}
\begin{proof}
\leanok
Apply the final assembly to the proved square reduction. Sample zero is
the entire boundary, whose squared distance from the fixed center is
$\ell^2/2>0$. The disk bridge identifies the whole body.
\end{proof}

\begin{proposition}[A nondisk would give a nonconstant geometric pentagon solution]
\label{prop:pentagon-counterexample}
\lean{UlamFloating.IsUnitSpeedBoundary.nonconstant_pentagon_solution_of_not_isDisk,UlamFloating.IsUnitSpeedBoundary.nonconstant_pentagon_solution_with_lift_of_not_isDisk}
\leanok
\uses{prop:constant-fifth-rigidity,thm:geometric-pentagon-ode,prop:pentagon-period,prop:chord-one-turn,prop:cyclic-angle-sum,prop:chord-angle-range}
If a constant one-fifth chord body is not a disk, its boundary admits a
reconstructing chord lift with a nonconstant $C^1$, $P$-periodic sampled
curve $X$. The coordinates lie in $(0,\pi)$, sum to $3\pi$, satisfy
$X_{i+1}(t)=X_i(t+P/5)$, and solve $X'=(2/\ell)F(X)$.
Every positive period $T$ of $X$, including a period smaller than $P$,
satisfies
\[
 T\geq\frac{\ell\pi}{2\sqrt3}.
\]
\end{proposition}
\begin{proof}
\leanok
Choose the positive one-turn lift. Its sampled curve cannot be constant,
since the constant-case theorem would make the body a disk. Apply the
geometric equation and the analytic period bound to this same curve.
The original boundary reconstruction is retained in the formal conclusion.
The stronger version also retains the $C^1$ lift, its positive derivative,
and its exact one-turn translation identity.
\end{proof}

\begin{definition}[The additive subgroup of periods]
\label{def:period-subgroup}
\lean{UlamFloating.periodSubgroup}
\leanok
For a function $f:G\to Y$ on any additive group, its period subgroup is
$\{T\in G: f(t+T)=f(t)\text{ for every }t\}$.
The target type $Y$ is arbitrary. Zero and negative periods are included.
\end{definition}

\begin{proposition}[A period gap gives a least positive period]
\label{prop:least-period-gap}
\lean{UlamFloating.exists_least_positive_period_of_lower_bound}
\leanok
\uses{def:period-subgroup}
In a linearly ordered Archimedean additive commutative group, suppose $f$
has a positive period and every positive period is at least $\varepsilon>0$.
Then $f$ has a least positive period $T$.
Continuity and a differential equation are not assumptions of this result.
\end{proposition}
\begin{proof}
\leanok
The period subgroup is nontrivial and disjoint from $(0,\varepsilon)$.
Apply mathlib's least-positive-element theorem for ordered Archimedean
subgroups.
\end{proof}

\begin{lemma}[Classification of all periods]
\label{lem:period-classification}
\lean{UlamFloating.mem_periodSubgroup,UlamFloating.periodSubgroup_eq_zmultiples_of_isLeast,UlamFloating.periodic_iff_zsmul_of_isLeast,UlamFloating.periodic_pos_iff_nat_mul_of_isLeast}
\leanok
\uses{def:period-subgroup,prop:least-period-gap}
Membership in the period subgroup is ordinary periodicity. If $T$ is
the least positive period, the subgroup is exactly $\mathbb ZT$.
For a real-valued time domain, a positive number $P$ is a period precisely
when $P=mT$ for a positive natural number $m$.
\end{lemma}
\begin{proof}
\leanok
A least positive element generates an Archimedean ordered subgroup.
Use the integer-multiple description of that generated subgroup.
Since $P,T>0$, the integer multiplying $T$ is positive and hence is a
positive natural number.
\end{proof}

\begin{lemma}[Cyclic shifts exclude divisible multiplicities]
\label{lem:cyclic-period-divisibility}
\lean{UlamFloating.eq_initial_of_finRotate_invariant,UlamFloating.constant_of_periodic_cyclic_shift,UlamFloating.cycle_length_not_dvd_period_multiple}
\leanok
\uses{lem:finite-cycle}
A function on a nonempty finite cycle invariant under one-step rotation
is constant. Consequently, if a coordinate motion $X:\mathbb R\to\mathbb R^N$
has fixed coordinate sum and satisfies
$X_{i+1}(t)=X_i(t+P/N)$, then $P/N$ being a period forces $X$ to be constant.
For nonconstant $X$, if $T$ is a period and $P=mT$ with $m\in\mathbb N$,
then $N$ does not divide $m$. No continuity or ODE is required here.
\end{lemma}
\begin{proof}
\leanok
Induct along the finite cycle to show all coordinates agree. The fixed
sum determines their common value. If $m=Nq$, then $P/N=qT$ is a period;
the preceding constancy conclusion contradicts nonconstancy.
\end{proof}

\begin{proposition}[Least periods of general cyclic coordinate motions]
\label{prop:least-cyclic-period}
\lean{UlamFloating.exists_least_period_cyclic_motion}
\leanok
\uses{prop:least-period-gap,lem:period-classification,lem:cyclic-period-divisibility}
For a nonconstant fixed-sum cyclic motion with positive traversal period
$P$ and positive lower bound $\varepsilon$ on periods, there are a least
positive period $T\geq\varepsilon$ and a positive integer $m$ with
$P=mT$ and $N\nmid m$.
\end{proposition}
\begin{proof}
\leanok
Choose the least positive period from the period gap. Classify the
traversal period as a positive multiple and apply the cyclic divisibility
obstruction. Apply the original lower bound to the least period itself.
\end{proof}

\begin{lemma}[Translation increments of a primitive with periodic derivative]
\label{lem:periodic-primitive-increment}
\lean{UlamFloating.sub_eq_of_periodic_deriv,UlamFloating.add_nat_mul_of_add_const}
\leanok
Let $f:\mathbb R\to E$ be differentiable, where $E$ is any real normed
space. If $f'$ is $T$-periodic, then $f(t+T)-f(t)$ is independent of $t$.
More generally, any identity $f(t+T)=f(t)+v$ implies
$f(t+nT)=f(t)+nv$ for every natural number $n$.
\end{lemma}
\begin{proof}
\leanok
The derivative of $f(t+T)-f(t)$ vanishes, so the difference is constant
by the mean value theorem. Iterate a fixed increment by induction.
\end{proof}

\begin{lemma}[Dividing a known total increment]
\label{lem:periodic-primitive-multiplicity}
\lean{UlamFloating.sub_eq_inv_smul_of_periodic_deriv,UlamFloating.sub_eq_div_of_periodic_deriv}
\leanok
\uses{lem:periodic-primitive-increment}
If $f'$ is $T$-periodic, $m>0$ is an integer, and $f(mT)-f(0)=A$,
then $f(t+T)-f(t)=m^{-1}A$ for every $t$. The result holds in any
real normed space; the scalar specialization gives $A/m$.
\end{lemma}
\begin{proof}
\leanok
The preceding lemma gives $f(mT)-f(0)=m(f(T)-f(0))$.
Multiply by $m^{-1}$ and use the independence of the initial time.
\end{proof}

\begin{lemma}[Geometric angular velocity from pentagon shape]
\label{lem:pentagon-angular-speed}
\lean{UlamFloating.IsUnitSpeedBoundary.fifth_chord_angle_speed}
\leanok
\uses{prop:boundary-tangent-angles,lem:five-angle-relations}
For the positive $C^1$ one-turn chord lift of a constant one-fifth chord
boundary, with sampled angle curve $X$, one has
\[
 \theta'(t)=-\frac{2}{\ell}\sin(X_1(t)+X_3(t)).
\]
Indices here are zero based, as in the formal definition.
\end{lemma}
\begin{proof}
\leanok
Apply the five tangent-angle identity and the geometric sine speed law.
Identify its two translated angle expressions with sampled coordinates.
\end{proof}

\begin{proposition}[Exact chord turning over a shape period]
\label{prop:pentagon-shape-turning}
\lean{UlamFloating.IsUnitSpeedBoundary.periodic_chord_angle_deriv_of_fifth_shape_period,UlamFloating.IsUnitSpeedBoundary.chord_angle_advance_of_fifth_shape_period}
\leanok
\uses{lem:pentagon-angular-speed,lem:periodic-primitive-multiplicity}
Every period $T$ of the actual sampled shape $X$ is a period of
$\theta'$. If $P=mT$ for a positive integer $m$, then
\[
 \theta(t+T)-\theta(t)=\frac{2\pi}{m}\qquad\text{for every }t.
\]
The lift is the original geometric chord lift, with its proved one-turn
identity over the boundary traversal.
\end{proposition}
\begin{proof}
\leanok
Periodicity of $X$ makes the angular-speed formula periodic. Apply the
primitive increment theorem with total increment $2\pi$ over $P=mT$.
\end{proof}

\begin{lemma}[Area-preserving maps and polygon moments]
\label{lem:polygon-map-moments}
\lean{UlamFloating.areaForm_rotation,UlamFloating.polygonSignedArea_map,UlamFloating.polygonFirstMoment_map,UlamFloating.polygonCentroid_map}
\leanok
\uses{def:moments}
Rotations preserve the oriented area form in any Euclidean plane.
For any planar linear map $L$ preserving that form and any finite
successor permutation, signed polygon area is unchanged, while signed
first moment and signed centroid are carried to their images under $L$.
The centroid identity also holds for the totalized zero-area definition.
\end{lemma}
\begin{proof}
\leanok
Expand the rotation into cosine and sine components and apply
$\cos^2 a+\sin^2 a=1$. For a general area-preserving map, apply its
linearity to the finite sums defining moments and centroid.
\end{proof}

\begin{lemma}[Recovering vertices from cyclic edge vectors]
\label{lem:cyclic-edge-reconstruction}
\lean{UlamFloating.eq_add_of_cyclic_edges_eq}
\leanok
\uses{lem:cyclic-period-divisibility}
For two polygons $p,q$ on the same nonempty finite cycle with values
in any additive commutative group, equality of all directed edge vectors
implies $q_i=p_i+(q_0-p_0)$ for every index $i$.
\end{lemma}
\begin{proof}
\leanok
The difference $q_i-p_i$ is unchanged by the cyclic successor.
Finite induction makes it equal to its value at index zero.
\end{proof}

\begin{proposition}[Recovering centered polygons from transformed edges]
\label{prop:centered-edge-reconstruction}
\lean{UlamFloating.sub_centroid_eq_of_edges_eq_map}
\leanok
\uses{lem:polygon-map-moments,lem:cyclic-edge-reconstruction,prop:moment-translation}
Suppose $p$ has nonzero signed area, $L$ preserves the oriented area
form, and every directed edge of $q$ is the image under $L$ of the
corresponding edge of $p$. Writing $C$ for signed centroid, one has
\[
 q_i-C(q)=L(p_i-C(p)).
\]
Convexity and regularity are not required.
\end{proposition}
\begin{proof}
\leanok
Edge reconstruction gives $q_i=L(p_i)+b$ for one translation $b$.
The centroid transformation identities give $C(q)=L(C(p))+b$.
Subtract to eliminate the translation.
\end{proof}

\begin{proposition}[Rotated carousel edges fix a common rotation center]
\label{prop:carousel-rotated-edges}
\lean{UlamFloating.IsCarousel.sub_centroid_eq_of_rotated_edges}
\leanok
\uses{prop:centered-edge-reconstruction,prop:invariants}
If two configurations of a nonzero-area carousel have edge vectors
related by the rotation $R_a$, their vertices are related by the same
rotation about the conserved signed centroid $c$:
$p_i(s)-c=R_a(p_i(t)-c)$.
\end{proposition}
\begin{proof}
\leanok
Apply centered edge reconstruction to the rotation, then use conservation
of the signed centroid to identify both centers with the centroid at time zero.
\end{proof}

\begin{proposition}[A chord-angle increment rotates the original boundary]
\label{prop:chord-boundary-rotation}
\lean{UlamFloating.rotation_of_angle_increment,UlamFloating.IsUnitSpeedBoundary.rotation_about_centroid_of_chord_angle_increment}
\leanok
\uses{prop:carousel-rotated-edges,prop:extraction,prop:polygon-area-nonzero}
For any constant rational chord increment $P/n$ with $n\geq3$,
suppose the original chord lift satisfies $\theta(t+T)=\theta(t)+a$.
Then, with $c$ the signed centroid of the sampled polygon at time zero,
\[
 \gamma(t+T)-c=R_a(\gamma(t)-c)\qquad\text{for every }t.
\]
The rotation center is derived from the geometric boundary data.
The underlying vector-rotation identity needs no differentiability.
\end{proposition}
\begin{proof}
\leanok
Composition of rotations turns the angle increment into rotation of
every directed chord. Apply this at all sampled vertices and use the
carousel result. Strict convexity supplies nonzero signed area. Sample
zero is the original boundary point at each time.
\end{proof}

\begin{proposition}[Least shape period of a hypothetical pentagon counterexample]
\label{prop:least-pentagon-period}
\lean{UlamFloating.IsUnitSpeedBoundary.least_pentagon_period_of_not_isDisk}
\leanok
\uses{prop:pentagon-counterexample,prop:least-cyclic-period,prop:pentagon-shape-turning,prop:chord-boundary-rotation}
A nondisk with the constant one-fifth chord property would give the
nonconstant geometric angle solution described above, with least positive
period $T$ and positive integer $m$ satisfying
\[
 P=mT,\qquad \gcd(m,5)=1,\qquad
 T\geq\frac{\ell\pi}{2\sqrt3}.
\]
The conclusion retains the chord reconstruction, convex angle ranges,
fixed sum, cyclic time shift, and actual differential equation.
It also gives $\theta(t+T)-\theta(t)=2\pi/m$ for every $t$.
For a single fixed center $c$, the original boundary satisfies
$\gamma(t+T)-c=R_{2\pi/m}(\gamma(t)-c)$ for every $t$.
\end{proposition}
\begin{proof}
\leanok
Apply the general least-period theorem to the extracted geometric motion.
Its analytic period bound is strictly positive because $\ell>0$.
The cyclic divisibility obstruction gives $5\nmid m$; primality of five
makes this equivalent to coprimality.
The strengthened extraction retains the lift data needed to apply the
exact shape-period turning theorem to the same angle and period.
The general chord-to-boundary rotation theorem identifies a fixed center
for the same time shift, with no assumption of regular sampled polygons.
\end{proof}

\begin{proposition}[Pentagon obstruction: unresolved input]
\label{prop:pentagon}
\lean{UlamFloating.pentagonObstructionStatement}
\notready
\uses{prop:extraction,prop:invariants,prop:least-pentagon-period,prop:constant-fifth-rigidity}
A strictly convex, once-traversed unit-speed boundary with constant chords
at increment $P/5$ must have regular sampled pentagons.
\end{proposition}
The constant geometric angle branch is proved, including the full regular
pentagon distance table. Excluding nonconstant geometric motions remains
unresolved in this project. The least positive shape period and the relation
$P=mT$ with $\gcd(m,5)=1$ are now proved, but these constraints do not
exclude every positive multiplicity. Possible
routes are a rigorous period obstruction for the pentagon flow or a direct
inequality excluding nonconstant motions subject to the one-curve shift
condition. Plots and approximate decimal periods are not proof inputs.
An interval-based approach needs exact certificates and a proved enclosure
theorem. Do not encode the desired obstruction as a geometric hypothesis.

\begin{proposition}[Fixed circle for a regular-polygon carousel]
\label{prop:regular-carousel}
\lean{UlamFloating.IsCarousel.fixed_circle_of_regular}
\leanok
\uses{prop:cyclic-centroid,prop:invariants,def:polygon-geometry}
A carousel whose polygons are regular and whose signed area is nonzero
at time zero has a single fixed circumcenter and positive squared
circumradius. Regularity is a hypothesis of this proposition.
\end{proposition}
\begin{proof}
\leanok
Area conservation keeps the signed area nonzero. The signed centroid is
the circumcenter at each time, and its conservation fixes that center.
The regular-polygon distance formula gives
$\ell^2=2r^2(1-\cos(2\pi/n))$ at every time. The coefficient is nonzero
because $\ell>0$, so the squared circumradius is also fixed.
\end{proof}

\begin{lemma}[Disk rigidity from regular rational samples]
\label{lem:regular-samples-disk}
\lean{UlamFloating.IsUnitSpeedBoundary.isDisk_of_regular_samples}
\leanok
\uses{prop:regular-carousel,prop:extraction,prop:polygon-area-nonzero,prop:disk}
For $n\geq3$, a constant-chord boundary with increment $P/n$ is a disk
if every sampled polygon is regular.
\end{lemma}
\begin{proof}
\leanok
Extract the carousel. Cyclic strict convexity supplies nonzero signed
area. The fixed-circle theorem applies, and sample zero traces the whole
boundary. Reconstruct the body from its spherical frontier.
\end{proof}

\begin{lemma}[One-fifth assembly with explicit obstruction premise]
\label{lem:fifth-assembly}
\lean{UlamFloating.fifth_rigidity_of_pentagon_obstruction}
\leanok
\uses{lem:regular-samples-disk,prop:pentagon}
The pentagon obstruction implies $\mathcal R(1/5)$.
\end{lemma}
\begin{proof}
\leanok
Apply the obstruction to the boundary data in $\mathcal P_{1/5}(K)$,
then use rigidity from regular rational samples with $n=5$.
The obstruction remains an explicit premise of this checked implication.
\end{proof}

\begin{theorem}[Density $1/5$: target, not yet proved]
\label{thm:fifth}
\lean{UlamFloating.fifthRigidityStatement}
\notready
\uses{prop:pentagon,lem:fifth-assembly}
$\mathcal R(1/5)$ holds.
\end{theorem}
The assembly after the obstruction is proved: conservation of the algebraic
centroid fixes the circumcenter, and equal edge lengths fix the radius.
No polygon-to-lamina centroid theorem is needed for this route.
The obstruction remains unproved and is not an axiom in the Lean library.

\section{Physical formulation and source audit}

\begin{definition}[Physical equilibrium and perimetral fraction]
\label{def:physical}
\lean{UlamFloating.waterlineCap,UlamFloating.IsEquilibriumCut,UlamFloating.FloatsWithPerimetralDensity}
\leanok
\uses{def:boundary,def:moments}
An equilibrium cut has unit normal $u$, offset $a$, cap area equal to
$\rho$ times body area, and cap-to-body centroid displacement parallel to $u$.
Floating at perimetral density $\sigma$ requires $0<\sigma<1$, boundary
data $(P,\gamma)$, and one $\rho\in(0,1)$ such that every pair
$\gamma(t),\gamma(t+\sigma P)$ lies on an equilibrium cut and every
oriented unit normal occurs. This definition does not assume constant
chord lengths, regular sampled polygons, or rigidity.
\end{definition}

\begin{proposition}[Hydrostatic reduction: statement formalized]
\label{prop:physical}
\lean{UlamFloating.hydrostaticReductionStatement}
\uses{def:physical,def:perimetral}
An intrinsic equilibrium condition in every orientation, with perimetral
density $\sigma$, implies $\mathcal P_\sigma(K)$.
\end{proposition}
The formal interface uses unit normals $u$, offsets $a$, and caps
$K\cap\{x:\ip ux\leq a\}$. Fix one submerged area fraction
$\rho\in(0,1)$ for all orientations. Equilibrium requires the cap centroid
minus the body centroid to be parallel to $u$. The waterline endpoints
must divide the boundary at increment $\sigma P$. Prove integrability,
positive cap area, unique offsets, endpoint regularity, and the necessary
cap-integral differentiation. A converse is a separate future theorem.
The physical rigidity statements also require the intrinsic boundary bridge.

\begin{lemma}[Physical assembly: conditional]
\label{lem:physical-assembly}
\lean{UlamFloating.physical_rigidity_of_geometric}
\leanok
\uses{def:physical,def:perimetral}
The hydrostatic reduction and the three geometric rigidity statements
together imply physical rigidity at the three densities. All four
statements are explicit premises of this assembly lemma.
\end{lemma}
\begin{proof}
\leanok
Apply hydrostatic reduction and split into the three values of $\sigma$.
\end{proof}

\begin{theorem}[Physical rigidity at the three densities: target]
\label{thm:physical-rigidity}
\lean{UlamFloating.physicalRigidityStatement}
\uses{prop:physical,thm:third,thm:quarter,thm:fifth,lem:physical-assembly}
If $\sigma\in\{1/3,1/4,1/5\}$ and a body floats at perimetral density
$\sigma$ in the preceding formal sense, then it is a closed disk of
positive radius. The boundary-data formulation is explicit; applying this
to a merely regular input boundary also requires arc-length normalization.
\end{theorem}

\paragraph{Source audit.}
Bracho--Montejano--Oliveros, \emph{Carousels, Zindler curves and the floating
body problem} (2004), Theorem 5, treats $1/3$ and $1/4$. Section 6 asserts
the pentagon period obstruction and refers elsewhere for details.
\href{https://www.matem.unam.mx/~luis/pdfs/volantin1.pdf}{Author-hosted PDF}.

Their \emph{A classification theorem for Zindler carrousels} (2001),
Section 4, labels period monotonicity Conjecture 4.1 and discusses computer
evidence. Its displayed inequality conflicts with the following numerical
discussion. Theorem 3.5 also excludes two exceptional cases by figures.
These steps need rigorous replacements or a verified later source; this
audit does not establish that no rigorous proof exists elsewhere.
\href{https://www.matem.unam.mx/~luis/pdfs/volantin2.pdf}{Author-hosted PDF}.

The 2026 sixth-density paper mentions numerical evidence in the earlier
fifth-density work, while the eighth-density paper summarizes the low
denominators as proved. Its own eighth-density theorem assumes central
symmetry, which is absent from our three targets.
\href{https://arxiv.org/abs/2604.10330}{Sixth-density paper};
\href{https://arxiv.org/abs/2607.08338}{Eighth-density paper}.

\paragraph{Implementation order.}
The triangle and square reductions and geometric rigidity at $1/3$ and $1/4$
are proved, as is the full convex carousel extraction theorem.
The analytic period bound, unit-speed metric length, and arc-length
normalization of regular boundary data are also proved.
The five-angle equation has an exact period lower bound from finite-cycle
energy, with positive time scaling explicit. Global chord-angle coordinates,
one-turn winding, periodic convex cyclic angle coordinates, and their
geometric five-angle differential equation are proved. The constant angle
branch gives both a disk and regular sampled pentagons. A nondisk would
therefore yield a nonconstant geometric solution, with the proved lower
bound on every positive shape period. The least positive shape period $T$
now exists, and $P=mT$ for a positive integer $m$ coprime to five.
A stronger geometric or period obstruction is still required to exclude
all such nonconstant motions; the present lower bound is insufficient.
Prove geometric centroid identification separately.
The final one-fifth assembly is proved conditional on the pentagon obstruction,
using signed moments to identify the circumcenter directly.
Audit and replace the pentagon input
before certifying $1/5$. Develop hydrostatics and intrinsic regularity as
separate modules. General $q/n$ sampling needs additional coprimality and
winding proofs. Keep metric lemmas polymorphic, local laws in real inner
product spaces, and finite-sum arguments independent of the denominator.
Normalize side length only through an explicit scaling equivalence if needed.
